%% 波-10 見えない中を調べる：超音波と振動の計測（第24回編集委員会 物理第3巻：組 波W5）。★1〜4、3段にしない。
%% 記号（計画書 C-1・B の波-10。この章の中だけの分け方）：
%%   音の速さ c（W-3）。この章では減衰係数 c を出さず、センサーは減衰比 ζ と固有角振動数 ω_n（固有振動数 f_n）で書く（W-4）。
%%   はね返す物までの距離 L（圧電定数 d_33 と分けるため d を距離に使わない）、往復の時間 t、波長 λ、振動数 f、角振動数 ω＝2πf。
%%   音響インピーダンス Z＝ρc（波-7）、音圧で数えた反射係数 r＝(Z_2−Z_1)/(Z_2+Z_1)。強さの割合は r² と 1−r²（透過の記号は作らない）。
%%   振動の変位 x、速度 v、加速度 a、振幅 A（W-2）。センサーの台（ケース）の変位 y、おもりの台に対するずれ z、ゆれの倍率 H（ω/ω_n と ζ で決まる）。
%%   圧電定数 d_33〔C/N〕、電荷 Q、力 F。回転の周波数 f_r、玉の数 n、玉の直径 D_b、玉の中心を結ぶ円の直径 D_p。ひずみ ε、ゲージ率 2。
%% 公式の持ち主：この章は中学・高校の公式を導かない（計画書 0 の1。こだまの計算は A5 の使い方）。\phkoushiki は置かない（K-8）。
%%   反射係数（A56）は波-7 ⑤、ドップラー（A31・A32）は波-8 ⑤⑦、1自由度の式・ζ は波-1 ⑦、強制振動の振幅は波-2 ⑤（6-13）を指す。
%% ほかの巻との約束：
%%   力-9 の道しるべ「この章の先に 波-10」：⑩でひずみゲージ（ゲージ率 約2、120 Ω、ブリッジ回路は電-5）で振動をはかる。
%%   数学 9-4（1310行）・9-12④：毎分1800回転＝30 Hz、1倍は回る部分の重さのかたより、2倍は軸のずれ、軸受けのきずの周波数 → ⑧。
%%   数学 9-6（1131行）「インパルスハンマーでたたいてインパルス応答をはかる（物理 波-10）」→ ⑨。
%%   波-9：空気の音の速さ 343 m/s、水 約1480 m/s、こだまは往復（波-9 のまちがえやすいところ）。
%% 数値はすべて Python で確かめた：checks/物理第3巻/波-10/calc_p3_10.py（プログラムの出力は phcode_sensor.py）。
%% 出典（第24回の追記：書くときに確かめる。11月5日、組 波W5。この環境では WebFetch が DNS で届かなかったので、
%%   WebSearch で一次資料を引く資料を2つ以上照らし合わせた）：
%%   ・コウモリの超音波：国内の小型コウモリ類の多くは 10〜120 kHz（国土技術政策総合研究所 ks0354024.pdf）。キクガシラコウモリのドップラーシフト補償
%%     （同資料、名古屋市の調査報告）。仲間の声を避けて周波数をずらす種（Bates ら J. Exp. Biol.）、ずらさないという報告（Götze ら Sci. Rep. 2016、
%%     背負わせたマイクの記録 PMC4262180）。しくみの全体は議論が続いている。〔確〕
%%   ・駐車支援の超音波センサー 40〜70 kHz（EE Times Japan 2020/04/24）、40 kHz の製品（村田製作所）。〔確〕
%%   ・魚群探知機 50 kHz（広い・深い）と 200 kHz（細かい・浅い）：Raymarine の取扱説明書、Louisiana Sportsman の解説。〔確〕
%%   ・鋼の縦波 約5900 m/s：三菱重工の超音波探傷の資料、Olympus の音速表（1020 鋼 5890 m/s）。Z＝7850×5900＝4.63×10^7 Pa·s/m（計算）。〔確〕
%%   ・体（軟らかい組織）1540 m/s、Z 1.63×10^6 Pa·s/m、減衰 約0.5 dB/(cm·MHz)、深く見るには振動数を下げる：Stanford RAD225 Class3、
%%     チェコ工科大学 ultrasound.pdf（2.5〜10 MHz）。〔確〕
%%   ・PZT の d_33 約300〜700 pC/N（IEEE の解説）、約600 pC/N（Ceramics Bulletin）。水晶の d_11 約2.3 pC/N（arXiv:2104.02211 で 2.24 pm/V）。〔確〕
%%   ・キュリー兄弟の圧電効果 1880年（IEEE の解説ほか）。ランジュバンの水晶の振動子 1917年、潜水艦のこだま 1918年（Physics Today 2022、
%%     BMUS の F. Duck の年表、IEEE Milestone の提案）。フェッセンデンのこだまの距離計の特許 1914年・氷山の発見を目的（EBSCO、IEEE の議論）。〔確〕
%%   ・水晶時計 32 768 Hz（2の15乗）：計算。
%%   ・圧電式の加速度計：固有振動数の約3分の1まで使う（そこで約＋12 %）、使える上限 5〜50 kHz（HBK（B&K）の解説、All About Circuits）。〔確〕
%%   ・玉軸受けのきずの周波数 (n/2)f_r(1∓(D_b/D_p)cos φ)、2つの和が n f_r：arXiv:1503.04444、sensemore の計算の解説、fabrico の解説。〔確〕
%%   ・道路トンネルの点検（近接目視と打音検査、5年に1回）：国土交通省 道路トンネル定期点検要領（平成26年6月、令和6年9月改訂）、
%%     国土交通省の報道資料 road01_hh_000364。〔確〕
%%   ・新幹線の車軸：中空車軸の超音波探傷システム（イプロスの保守サービスの案内）、車軸の疲労き裂の研究（J-STAGE）。検査の周期の数値は書いていない。〔確〕
%%   ・スマホの加速度センサーが MEMS であること（Lumenci の解説ほか）。固有振動数の数値は本文に書いていない。
%%   ・探傷の振動数 1 MHz 以上、鋼の中の 5 MHz の縦波の波長 約1.2 mm（日本製鉄テクノロジー ut.pdf）、溶接部の探触子 おもに 2〜5 MHz（Screening Eagle の解説、イプロスの製品仕様）。〔確〕
%%   〔要確認〕のまま：なし（確かめられなかった数値は本文から外した）。
\makeatletter
\tikzset{rxk lab/.style={font=\sffamily\footnotesize, inner sep=1pt},
  rxk arr/.style={-{Stealth[length=3.6pt, width=3.2pt]}, line width=1pt},
  rxk big/.style={-{Stealth[length=4.2pt, width=3.8pt]}, line width=1.4pt},
  rxk dim/.style={{Stealth[length=3pt]}-{Stealth[length=3pt]}, line width=0.5pt, bkGray}}
% ばね（たて）  \rxkspring{(x,y下)}{長さ}
\newcommand\rxkspring[2]{\begin{scope}[shift={#1}]
  \draw[bkInk, line width=0.8pt] (0,0) -- (0,1.5) -- (-1.8,2.5) \foreach \k in {1,...,4} {-- (1.8,{2.5+(\k-0.5)*(#2-5)/4}) -- (-1.8,{2.5+\k*(#2-5)/4})} -- (0,{#2-1.5}) -- (0,#2);
\end{scope}}
% ダンパー（たて）  \rxkdamper{(x,y下)}{長さ}
\newcommand\rxkdamper[2]{\begin{scope}[shift={#1}]
  \draw[bkInk, line width=0.8pt] (0,0) -- (0,{#2*0.35}) (-2.2,{#2*0.35}) -- (2.2,{#2*0.35}) (-2.2,{#2*0.35}) -- (-2.2,{#2*0.72}) (2.2,{#2*0.35}) -- (2.2,{#2*0.72});
  \draw[bkInk, line width=0.8pt] (-1.5,{#2*0.55}) -- (1.5,{#2*0.55}) (0,{#2*0.55}) -- (0,#2);
\end{scope}}
\makeatother

\phchapter{10}{見えない中を調べる：\\超音波と振動の計測}{コウモリは、暗やみでどうやってものの位置を知る？}{\bklevelrange{1}{4}}{見えない中を調べる：超音波と振動の計測}
\bkmichishirube{9-4,9-6,6-13,力-9,波-1,波-2,波-7,波-8,波-9}{波-11,波-14}

\begin{goal}
\begin{itemize}[label={\color{bkTheme}●}, leftmargin=1.6\zw, itemsep=0pt, topsep=1pt]
\item \bklevel{1}\enspace こだまの時間から、コウモリ・ソナー・超音波探傷の距離を求められる。
\item \bklevel{3}\enspace 振動の変位・速度・加速度を換算し、回転機械のスペクトルから故障の見当をつけられる。
\item \bklevel{4}\enspace 1自由度の模型で、加速度計と変位計のしくみを説明し、感度のずれを計算できる。
\end{itemize}
\end{goal}

\begin{hatena}
\noindent\begin{minipage}[c]{0.56\linewidth}\raggedright
夕方の空を、コウモリがすばやく飛びまわって、小さな虫をつかまえています。暗くて、目ではほとんど見えないはずです。

コウモリは、どうやって、虫や木の枝がどこにあるのかを知っているのでしょう。
\end{minipage}\hfill
\begin{minipage}[c]{0.40\linewidth}\centering
\begin{tikzpicture}[x=0.8mm, y=0.8mm]
  % コウモリ
  \filldraw[fill=bkInk!75, draw=bkInk, line width=0.5pt] (8,20) .. controls (4,26) and (0,25) .. (-2,22) .. controls (1,21) and (2,19) .. (2,17) .. controls (4,18) and (6,17) .. (7,16) -- (10,18) -- (13,16) .. controls (14,17) and (16,18) .. (18,17) .. controls (18,19) and (19,21) .. (22,22) .. controls (20,25) and (16,26) .. (12,20) -- cycle;
  \filldraw[fill=bkInk!75, draw=bkInk] (10,19.5) circle (2.3);
  \fill[bkInk] (8.7,21.3) -- (9.3,23.6) -- (10,21.5) -- (10.7,23.6) -- (11.3,21.3) -- cycle;
  % 出す音（実線）
  \foreach \rr in {6,9,12} {\draw[bkBlue!80!black, line width=0.8pt] ([shift={(-25:\rr)}]23,19) arc (-25:25:\rr);}
  % 虫
  \filldraw[fill=bkOrange!60, draw=bkInk, line width=0.5pt] (52,20) ellipse (2.6 and 1.4);
  \draw[bkInk, line width=0.5pt] (51,21) -- (49,24) (53,21) -- (55,24);
  % 返る音（破線）
  \foreach \rr in {6,10} {\draw[bkOrange!85!black, line width=0.8pt, dashed] ([shift={(155:\rr)}]49,19) arc (155:205:\rr);}
  \node[rxk lab, anchor=north, text=bkBlue!70!black] at (26,9) {出す音（実線）};
  \node[rxk lab, anchor=north, text=bkOrange!70!black] at (46,3) {こだま（破線）};
\end{tikzpicture}
\end{minipage}
\end{hatena}

\noindent 波-9 では、耳に聞こえる音の速さと大きさを数にしました。この章では、耳に聞こえない高い音（\textbf{超音波}）と、機械や建物の\textbf{振動}を「はかる道具」として使います。コウモリのように、音を出してこだまを聞けば、見えない所にある物の場所がわかります。鉄の板の中の傷、おなかの中の赤ちゃん、海の中の魚の群れも、同じしくみで見えるのです。後半では、ゆれをはかるセンサーのしくみを、第1部のばねとおもりの模型で考えます。

%% ---------------- やってみよう・思い出そう ----------------
\begin{yattemiyou}[スマホで洗濯機のゆれをはかろう]
\textbf{用意するもの}\par
スマートフォン（加速度センサーの値をグラフで表示するアプリを入れる。大人に相談する）、洗濯機（または、ゆれる物）
\begin{enumerate}[label={\bkcircled[bkGreen]{\arabic*}}, leftmargin=2.4\zw, itemsep=1pt, topsep=2pt]
\item スマホを洗濯機のふたの上に平らに置き、テープで軽くとめる。動いている洗濯機の中には手を入れない。
\item 脱水のときの、上下の加速度のグラフを記録する。
\item グラフの山が1秒に何回あるかを数え、振動数（Hz）を読む。その振動数を60倍して、脱水槽の1分間の回転数と比べる。
\end{enumerate}
\end{yattemiyou}

\noindent スマホの中には、小さな\textbf{加速度センサー}が入っています。画面を横にすると表示が回るのは、このセンサーが重力の向きを感じているからです。センサーの中には、ばねで支えられた小さなおもりがあり、そのずれで加速度をはかっています。なぜ、ばねとおもりで加速度がはかれるのかは、⑦で調べます。

\begin{omoidasou}
\bkrecall{波-9}{空気の音の速さ $c\approx\qty{343}{m/s}$（\qty{20}{\celsius}）、水 約\qty{1480}{m/s}（海水 約\qty{1500}{m/s}）。こだまは往復なので、距離は（速さ×時間）÷2。}
\bkrecall{波-7}{音響インピーダンス $Z=\rho c$。境目ではね返る音圧の割合（反射係数）は $r=(Z_2-Z_1)/(Z_2+Z_1)$（⑤）。はね返る強さの割合は $r^2$。}
\bkrecall{波-8}{動く物ではね返った波の振動数のずれ $\Delta f\approx2vf/c$（⑦）。}
\bkrecall{波-1}{ばねとおもりの固有角振動数 $\omega_{\mathrm n}=\sqrt{k/m}$、$f_{\mathrm n}=\omega_{\mathrm n}/(2\pi)$。ダンパーの効き方は減衰比 $\zeta$（⑦、6-7⑤）。}
\bkrecall{波-2}{ゆすり続けると、ゆれの大きさはゆする振動数で変わる。$\omega=\omega_{\mathrm n}$ の近くで大きくなる（共振。6-13）。}
\bkrecall{9-4}{FFT で、ゆれを振動数ごとの成分（スペクトル）に分けられる（⑬）。}
\end{omoidasou}

%% ---------------- 物理の道具 ----------------
\Needspace{16\baselineskip}
\bksection{doguu}{物理の道具}

\bksubsection{1}{こだまで距離をはかる}

\leavevmode\label{ph:波-10:sec1}山で「ヤッホー」と叫ぶと、少しおくれて、向かいの山から「ヤッホー」と返ってきます。波-9 で学んだように、こだまが返ってくるまでの時間 $t$ の間に、音は物まで行って、もどってきます。物までの距離を $L$ とすると、音が進んだ道のりは $2L$ なので
\pheq{L=\frac{c\,t}{2}}{m ＝ m/s × s}
\label{ph:波-10:kodama}です。速さ $c$ がわかっていれば、時間をはかるだけで、見えない所にある物までの距離がわかります。

\textbf{はてな？の答え}：コウモリは、口や鼻から、人には聞こえない高い音（超音波。日本の小さなコウモリの多くは、種類によって約\qty{10}{kHz}〜\qty{120}{kHz}）を短く区切って出し、虫や枝からはね返ってくるこだまを聞いています。こだまが返るまでの時間で\textbf{距離}を、左と右の耳に届く時間や強さのわずかな差で\textbf{向き}を知ります。さらに、こだまの高さのずれ（ドップラー効果。波-8）で、相手が近づいているか遠ざかっているかもわかります。こだまが6.0 ms（1000分の6秒）後に返ってきたら、虫までの距離は
\[ L=\frac{343\times0.0060}{2}\approx\qty{1.0}{m} \]
です（\bkreidairef{1}）。このように、自分で出した音のこだまで周りを知るやり方を\textbf{エコーロケーション}（反響定位）といいます。\bkidx{えこーろけーしょん}{エコーロケーション}

\begin{bkfigure}
\begin{tikzpicture}[x=1mm, y=1mm]
  \draw[bkInk, line width=0.7pt, -{Stealth[length=3pt]}] (0,0) -- (120,0) node[rxk lab, anchor=west] {時間};
  % 出した音
  \draw[bkBlue!80!black, line width=1pt, domain=4:12, samples=80] plot (\x, {9*sin(180*(\x-4))*exp(-0.25*(\x-8)^2)});
  \node[rxk lab, anchor=south, text=bkBlue!70!black] at (8,11) {出した音};
  % こだま
  \draw[bkOrange!85!black, line width=1pt, domain=84:92, samples=80] plot (\x, {4*sin(180*(\x-84))*exp(-0.25*(\x-88)^2)});
  \node[rxk lab, anchor=south, text=bkOrange!70!black] at (88,6) {こだま（小さい）};
  \draw[rxk dim] (8,-6) -- (88,-6);
  \node[rxk lab, anchor=north] at (48,-7) {往復の時間 $t$（\qty{6.0}{ms}なら $L\approx\qty{1.0}{m}$）};
  \draw[bkGray, line width=0.4pt, densely dotted] (8,-8) -- (8,-1) (88,-8) -- (88,-1);
\end{tikzpicture}
\bkfigcaption{出した音と、こだまの時間の差。こだまは、広がったり、はね返る物が小さかったりして、出した音より小さい。}
\label{ph:波-10:fig-kodama}
\end{bkfigure}

\begin{phmijika}[車の後ろの障害物センサー]
車のバンパーについている丸いセンサーは、約\qty{40}{kHz}の超音波を出し、壁や人からのこだまで距離をはかっています。こだまが\qty{2.9}{ms}後に返れば、約\qty{0.50}{m}です（\bkreidairef{2}）。近づくと「ピッピッ」の間が短くなるように作られています。
\end{phmijika}

\begin{phbaisuru}{こだまの時間 $t$}
距離も2倍になります。$L=ct/2$ は $t$ に比例するからです。音の速さ $c$ をまちがえると、距離も同じ割合でずれます。気温が\qty{0}{\celsius}（\qty{331}{m/s}）と\qty{30}{\celsius}（\qty{349}{m/s}）では約5\,\% 違うので、正確な距離計は、気温をはかって音の速さを直しています（波-9 ⑤）。
\end{phbaisuru}

%% ---------------------------------------------------------------
\bksubsection{2}{なぜ超音波を使うのか}

\leavevmode\label{ph:波-10:sec2}こだまで物を見分けるには、なぜ、聞こえる音ではなく、超音波を使うのでしょうか。理由は\textbf{波長}にあります。

波は、波長よりずっと小さな物に当たっても、ほとんどはね返らず、回りこんで進んでしまいます（回折。波-8）。だから、こだまで見分けられる物の大きさは、だいたい波長くらいまでです。空気の中で\qty{1}{kHz}の音の波長は約\qty{34}{cm}ですが、\qty{100}{kHz}の超音波なら約\qty{3.4}{mm}です。コウモリが、小さな虫を見つけられるのは、波長の短い超音波を使っているからです。

\noindent\textbf{高すぎると、深くまで届かない}\par\nobreak
では、振動数は高ければ高いほどよいのでしょうか。そうではありません。体や物の中を進む超音波は、少しずつ熱に変わって弱まっていきます（減衰。波-7 ⑦）。体の軟らかい所では、強さが\qty{1}{cm}進むごとに、約 0.5\,dB×（振動数の MHz の数）だけ弱まります。振動数が高いほど、早く弱まるのです。こだまは行って帰るので、深さの2倍の道のりで弱まります。

\begin{bkfigure}
\begin{tikzpicture}
\begin{axis}[bkplot, width=104mm, height=56mm, xmin=0, xmax=1.8, ymin=0, ymax=68, xtick={0.2,0.4,0.6,0.8,1.0,1.2,1.4,1.6}, xticklabels={0.2,0.4,0.6,0.8,1.0,1.2,1.4,1.6}, ytick={10,20,30,40,50,60}, xlabel={波長（mm）：見分けられる大きさの目安}, ylabel={届く深さ（cm）}, axis x line*=bottom, axis y line*=left, every axis x label/.style={at={(axis description cs:0.5,-0.15)}, anchor=north, font=\sffamily\small}]
  \addplot[bkGray, line width=0.8pt, densely dotted, domain=0:1.62] {60/1.54*x};
  \fill[bkBlue] (axis cs:1.54,60) circle (2pt); \node[rxk lab, anchor=west] at (axis cs:1.57,58) {1 MHz};
  \fill[bkBlue] (axis cs:0.77,30) circle (2pt); \node[rxk lab, anchor=west] at (axis cs:0.80,28) {2 MHz};
  \fill[bkBlue] (axis cs:0.44,17.1) circle (2pt); \node[rxk lab, anchor=west] at (axis cs:0.47,15) {3.5 MHz};
  \fill[bkBlue] (axis cs:0.308,12) circle (2pt); \node[rxk lab, anchor=south east] at (axis cs:0.29,13) {5 MHz};
  \fill[bkBlue] (axis cs:0.154,6) circle (2pt); \node[rxk lab, anchor=north west] at (axis cs:0.18,5) {10 MHz};
  \node[rxk lab, anchor=west, align=left] at (axis cs:0.95,16) {振動数を上げると、左下へ：\\細かく見えるが、浅い所しか見えない};
\end{axis}
\end{tikzpicture}
\bkfigcaption{体の中の超音波の振動数と、見分けられる大きさ（波長）・届く深さ。どちらも振動数に反比例するので、点は1本の直線（点線）の上に並ぶ。届く深さは、往復で60\,dB 弱まるまでの深さ。体の中の音の速さ\qty{1540}{m/s}、減衰を 0.5\,dB/(cm·MHz) として計算。}
\label{ph:波-10:fig-tradeoff}
\end{bkfigure}

細かく見たいなら振動数を上げ、深くまで見たいなら下げる。この「見分けやすさ」と「届く深さ」の取り引きで、使う振動数が決まります（図\ref{ph:波-10:fig-tradeoff}、\bkreidairef{3}）。おなかの中を見る検査では約\qty{2}{MHz}〜\qty{5}{MHz}、体の表面に近い血管や甲状腺を見るときは約\qty{10}{MHz}以上が使われます。

\begin{phmodel}[こだまで物を見分ける]
\phignore{波が広がって弱まる分（どの振動数でも同じ割合で効く）。体の中の場所による音の速さの違い。}
\phassume{音はまっすぐ進み、速さは一定。見分けられる大きさは、だいたい波長くらい。減衰の割合は振動数に比例する。}
\phlimit{骨や空気（肺・腸のガス）のように、$Z$ が大きく違う所（ほとんどはね返って、その先が見えない）。音の速さが場所で大きく変わる所（距離がずれる）。}
\end{phmodel}

%% ---------------------------------------------------------------
\bksubsection{3}{鉄の中の傷を見つける：パルス反射法}

\leavevmode\label{ph:波-10:sec3}橋の鉄骨や、発電所の配管、新幹線の車軸などの中に、目に見えない割れ（傷）があると、やがて大きな事故につながります。部品をこわさずに中を調べる方法を\textbf{非破壊検査}といい、\bkidx{ひはかいけんさ}{非破壊検査}その代表が\textbf{超音波探傷}です。\bkidx{ちょうおんぱたんしょう}{超音波探傷}


\textbf{探触子}（プローブ）と呼ばれる器具を、板の表面に当てます。\bkidx{たんしょくし}{探触子}探触子は、1 MHz 以上（鋼の溶接部では、おもに2〜5 MHz）の超音波を、ほんの数回だけゆれる短い音（\textbf{パルス}）にして出し、すぐに「聞く」側にかわって、こだまを待ちます。板の中に傷がなければ、こだまは板の底（底面）からだけ返ってきます。途中に傷があると、底面より早く、傷からのこだまが返ってきます（図\ref{ph:波-10:fig-tansho}）。これを\textbf{パルス反射法}といいます。\bkidx{ぱるすはんしゃほう}{パルス反射法}

\begin{bkfigure}
\begin{tikzpicture}[x=1mm, y=1mm]
  % 鋼の板
  \fill[bkGray!25] (0,0) rectangle (58,30);
  \draw[bkInk, line width=0.8pt] (0,0) rectangle (58,30);
  \node[rxk lab, anchor=north] at (29,-1) {鋼の板（厚さ\qty{100}{mm}）};
  % 傷
  \draw[bkRed!75!black, line width=1.8pt] (24,18) -- (36,16);
  \node[rxk lab, anchor=west, text=bkRed!70!black] at (37,17) {傷};
  % 探触子
  \filldraw[fill=bkBlue!35, draw=bkInk, line width=0.7pt] (24,30) rectangle (34,37);
  \node[rxk lab, anchor=south] at (29,37.5) {探触子（出す・聞く）};
  % 超音波の道
  \draw[rxk arr, bkBlue!80!black] (27,29.5) -- (27,18.6);
  \draw[rxk arr, bkOrange!85!black, dashed] (31,17.4) -- (31,29.5);
  \draw[rxk arr, bkBlue!80!black] (21,29.5) -- (21,0.6);
  \draw[rxk arr, bkOrange!85!black, dashed] (18,0.6) -- (18,29.5);
  % 画面
  \draw[bkInk, line width=0.7pt, -{Stealth[length=3pt]}] (68,2) -- (124,2) node[rxk lab, anchor=north east] at (124,1) {時間};
  \draw[bkInk, line width=0.7pt] (68,2) -- (68,32);
  \draw[bkBlue!80!black, line width=1pt] (70,2) -- (70,28) -- (72,2);
  \draw[bkRed!75!black, line width=1pt] (88,2) -- (89,15) -- (90,2);
  \draw[bkInk, line width=1pt] (103,2) -- (104,22) -- (105,2);
  \node[rxk lab, anchor=south] at (71,28.5) {T};
  \node[rxk lab, anchor=south] at (89,15.5) {F};
  \node[rxk lab, anchor=south] at (104,22.5) {B};
  \draw[rxk dim] (70,-3) -- (89,-3);
  \node[rxk lab, anchor=north] at (79.5,-4) {\qty{20}{\micro s}};
  \draw[rxk dim] (70,-10) -- (104,-10);
  \node[rxk lab, anchor=north] at (87,-11) {\qty{34}{\micro s}};
\end{tikzpicture}
\bkfigcaption{パルス反射法。左：探触子から出た超音波は、傷（実線→破線）と、板の底（底面）ではね返る。右：画面の波形。T は出した音、F は傷のこだま、B は底面のこだま。時間から深さがわかる。}
\label{ph:波-10:fig-tansho}
\end{bkfigure}

鋼の中の縦波の速さは約\qty{5900}{m/s}です（波-9 ④の「広い固体」の値。探傷器のメーカーの表でも約\qty{5900}{m/s}）。傷のこだまが\qty{20}{\micro s}（100万分の20秒）後に返れば、傷の深さは
\[ L=\frac{5900\times20\times10^{-6}}{2}\approx\qty{0.059}{m}=\qty{59}{mm} \]
です（\bkreidairef{4}）。画面の上のこだまの位置が、そのまま深さのものさしになるのです。

\noindent\textbf{なぜ、傷ではね返るのか}\par\nobreak
傷の正体は、鋼の中にできた、とても薄いすき間です。すき間の中は空気（または、からっぽ）です。鋼の音響インピーダンスは約 $4.6\times10^7$\,Pa·s/m、空気は約413 Pa·s/m で、10万倍ほど違います。音圧で数えた反射係数（波-7 ④の $r_{\text{押}}$。この章では $r$ と書く）は
\[ r=\frac{Z_{\text{空気}}-Z_{\text{鋼}}}{Z_{\text{空気}}+Z_{\text{鋼}}}\approx-1 \]
で、音はほとんど全部はね返ります（$-$ は、音圧の山と谷が入れかわってはね返ること）。髪の毛より薄い割れでも、超音波にとっては「壁」なのです。

\noindent\textbf{探触子と板の間には、水や油をぬる}\par\nobreak
探触子と板の間に空気のすき間があると、同じ理由で、超音波は板に入る前にはね返ってしまいます。そこで、探触子と板の間には、水・油・ゼリーなどをぬります（波-7 のはてな？と同じ）。部品を水の中に沈めて、水を通して超音波を入れる方法（水浸法）もあります。水から鋼へ入るときの反射係数は
\[ r=\frac{Z_{\text{鋼}}-Z_{\text{水}}}{Z_{\text{鋼}}+Z_{\text{水}}}=\frac{4.63\times10^7-1.48\times10^6}{4.63\times10^7+1.48\times10^6}\approx0.94 \]
で、強さの $r^2\approx0.88$（88\,\%）がはね返り、鋼に入るのは約12\,\% です（\bkreidairef{5}）。それでも、空気のすき間よりは、ずっと多く入ります。

\begin{phyomi}[$L=ct/2$ と反射係数]
「こだまの時間で深さが、こだまの大きさで、はね返した物の性質がわかる」。インピーダンスの差が大きい境目ほど、大きなこだまが返ります。傷が大きいほど、はね返す面積が広く、こだまも大きくなります。
\end{phyomi}

%% ---------------------------------------------------------------
\bksubsection{4}{厚さをはかる・海の中を見る}

\leavevmode\label{ph:波-10:sec4}\noindent\textbf{超音波厚さ計}\par\nobreak
古いタンクや配管は、内側がさびて、外から見えないまま薄くなっていくことがあります。超音波厚さ計は、外側から探触子を当てて、底面（内側の面）からのこだまの時間で、厚さをはかります。厚さ\qty{10}{mm}の鋼なら、往復の時間は $2\times0.010\div5900\approx\qty{3.4}{\micro s}$ です。これが\qty{3.1}{\micro s}に減っていたら、厚さは約\qty{9.1}{mm}で、\qty{1}{mm}ほど薄くなっています（\bkreidairef{6}）。100万分の1秒より細かく時間をはかる電子回路のおかげで、\qty{0.1}{mm}くらいの差がわかります。

\noindent\textbf{ソナーと魚群探知機}\par\nobreak
海の中では、光も電波もすぐに弱まって、遠くまで届きません。しかし、音は水の中でよく伝わります。船の底から真下に超音波を出し、海の底や魚の群れからのこだまで、深さをはかるのが\textbf{魚群探知機}や\textbf{ソナー}です。\bkidx{ぎょぐんたんちき}{魚群探知機}海水の音の速さは約\qty{1500}{m/s}なので、こだまが0.20秒後に返れば、深さは $1500\times0.20\div2=\qty{150}{m}$ です（\bkreidairef{7}）。魚群探知機は、深くまで届く約\qty{50}{kHz}（波長\qty{30}{mm}）と、細かく見分けられる約\qty{200}{kHz}（波長\qty{7.5}{mm}）を使い分けています（②と同じ取り引き）。

魚の群れが船に近づいたり遠ざかったりしていると、こだまの振動数が少しずれます。波-8 ⑦のドップラー効果の式 $\Delta f\approx2vf/c$ で、群れの速さがわかります（\bkreidairef{8}）。

\begin{column}{圧電とソナーの始まり}
1880年、フランスのピエール・キュリーと兄のジャックは、水晶などの結晶を押すと、表面に電気が出ることを発見しました（⑤の圧電効果）。キュリーは、のちにマリー・キュリーと結婚し、放射能の研究で知られる人です。

1912年、大きな客船タイタニック号が氷山にぶつかって沈んだころから、水の中の見えない物を音で見つける研究が進みました。1914年には、カナダ生まれのフェッセンデンが、こだまで距離をはかる装置の特許をとっています。第一次世界大戦中の1917年、フランスのランジュバンは、水晶に電気をかけて超音波を出す装置を作り、翌年には潜水艦からのこだまを聞くことに成功しました。これが、ソナーや魚群探知機、そして体の中を見る超音波の検査の始まりです。
\end{column}

%% ---------------------------------------------------------------
\bksubsection{5}{超音波を出して聞く：圧電素子}

\leavevmode\label{ph:波-10:sec5}探触子の中で、超音波を出したり聞いたりしているのは、\textbf{圧電素子}という小さな板です。\bkidx{あつでんそし}{圧電素子}水晶や、チタン酸ジルコン酸鉛（PZT）というセラミックスの板を押し縮めると、板の上と下の面に、＋と－の電気がたまります。これを\textbf{圧電効果}といいます。\bkidx{あつでんこうか}{圧電効果}ガスコンロやライターの「カチッ」で火花が出るのは、ばねで圧電素子を強くたたいて、高い電圧を作っているからです。

\begin{bkfigure}
\begin{tikzpicture}[x=1mm, y=1mm]
  % 押す
  \filldraw[fill=bkOrange!30, draw=bkInk, line width=0.7pt] (10,0) rectangle (40,8);
  \foreach \x in {13,19,25,31,37} {\node[rxk lab, text=bkRed!75!black] at (\x,6) {＋}; \node[rxk lab, text=bkBlue!75!black] at (\x,2) {－};}
  \draw[rxk big] (25,18) -- (25,9);
  \draw[rxk big] (25,-10) -- (25,-1);
  \node[rxk lab, anchor=west] at (26.5,15) {力 $F$};
  \node[rxk lab, anchor=north] at (25,-11) {押すと、電気が出る（聞く）};
  % 電気をかける
  \filldraw[fill=bkOrange!30, draw=bkInk, line width=0.7pt] (78,-1) rectangle (108,9);
  \draw[bkGray, line width=0.6pt, dashed] (80,0) rectangle (106,8);
  \draw[bkInk, line width=0.7pt] (93,9) -- (93,14) -- (115,14) -- (115,-6) -- (93,-6) -- (93,-1);
  \draw[bkInk, line width=0.7pt] (113,5) -- (117,5) (114,3) -- (116,3);
  \node[rxk lab, anchor=west] at (117.5,4) {電気};
  \draw[rxk arr] (93,15.5) -- (93,20);
  \draw[rxk arr] (93,-7.5) -- (93,-11);
  \node[rxk lab, anchor=north] at (93,-12) {電気をかけると、のび縮みする（出す）};
\end{tikzpicture}
\bkfigcaption{圧電素子。左：押し縮めると、上下の面に電気がたまる。右：電気をかけると、厚さがのび縮みする（破線はもとの大きさ）。}
\label{ph:波-10:fig-atsuden}
\end{bkfigure}

たまる電気の量（電荷）$Q$ は、押す力 $F$ に比例します。
\pheq{Q=d_{33}\,F}{C ＝ C/N × N}
比例定数 $d_{33}$ を\textbf{圧電定数}といいます（33 は「3の向きに押して、3の向きの面に出る」という意味の番号）。PZT では約\qty{400e-12}{C/N}（400 pC/N。材料によって300〜700 pC/N くらい）です。水晶は、押す向きと電気の出る向きの組み合わせが違い（$d_{11}$）、約 2.3 pC/N と小さな値です。\qty{10}{N}で押すと、PZT では $400\times10^{-12}\times10=\qty{4e-9}{C}$ がたまります（\bkreidairef{9}）。とても小さな電気ですが、電子回路で大きくしてはかれます。

反対に、圧電素子に電気をかけると、厚さがわずかにのび縮みします（逆圧電効果）。電気を\qty{5}{MHz}で＋と－に切りかえれば、板は\qty{5}{MHz}でふるえ、超音波を出します。1つの素子で「出す」と「聞く」の両方ができるのです。水晶時計の中では、水晶が1秒に\num{32768}回（2の15乗回）ふるえていて、それを数えて時間を決めています。電気の側のしくみ（コンデンサーとしての圧電素子）は、物理シリーズの『電磁気学』（電-3）で学びます。

%% ---------------------------------------------------------------
\bksubsection{6}{ゆれの3つの数：変位・速度・加速度}

\leavevmode\label{ph:波-10:sec6}ここからは、機械や建物の\textbf{振動}をはかります。振動の大きさは、変位（どれだけずれたか）・速度・加速度の3つの数で表せます。振動数 $f$、振幅 $A$ の正弦波のゆれ $x=A\sin\omega t$（$\omega=2\pi f$）なら、5-15 の微分で
\[ v=\frac{\dd x}{\dd t}=A\omega\cos\omega t,\qquad a=\frac{\dd v}{\dd t}=-A\omega^2\sin\omega t \]
です。それぞれの大きさ（振幅）は
\pheq{v_{\max}=A\omega=2\pi fA,\qquad a_{\max}=A\omega^2=(2\pi f)^2A}{m/s ＝ 1/s × m、m/s$^2$ ＝ 1/s$^2$ × m}
\label{ph:波-10:vxa}になります。

\begin{bkfigure}
\begin{tikzpicture}
\begin{axis}[bkplot, width=108mm, height=50mm, xmin=0, xmax=0.215, ymin=-1.3, ymax=1.3, xtick={0.05,0.1,0.15,0.2}, xticklabels={0.05,0.10,0.15,0.20}, ytick={-1,0,1}, xlabel={$t$（s）}, ylabel={}, axis x line*=bottom, axis y line*=left]
  \draw[bkGray, line width=0.5pt] (axis cs:0,0) -- (axis cs:0.21,0);
  \addplot[bkcurve, domain=0:0.2, samples=120] {sin(deg(2*pi*10*x))};
  \addplot[bkOrange, line width=1.5pt, dashed, domain=0:0.2, samples=120] {cos(deg(2*pi*10*x))};
  \addplot[bkGreen!60!black, line width=1.5pt, densely dotted, domain=0:0.2, samples=120] {-sin(deg(2*pi*10*x))};
\end{axis}
\end{tikzpicture}
\bkfigcaption{\qty{10}{Hz}・振幅\qty{1}{mm}のゆれ。実線は変位 $x$（振幅\qty{1}{mm}）、破線は速度 $v$（振幅\qty{63}{mm/s}）、点線は加速度 $a$（振幅\qty{3.9}{m/s^2}）。たての目もりは、それぞれの振幅を1とした。速度は変位より4分の1周期早く、加速度は変位と逆向き。}
\label{ph:波-10:fig-vxa}
\end{bkfigure}

\begin{phbaisuru}{振動数 $f$}
変位の振幅が同じでも、速度は2倍、加速度は4倍になります。\qty{10}{Hz}・\qty{1}{mm}なら速度 約\qty{63}{mm/s}・加速度 約\qty{3.9}{m/s^2}ですが（\bkreidairef{10}）、\qty{20}{Hz}なら\qty{126}{mm/s}・\qty{15.8}{m/s^2}（重力の加速度の1.6倍）です。速くゆれるものほど、わずかな変位でも大きな力（＝質量×加速度）がかかります。
\end{phbaisuru}

だから、はかるものは振動数で選びます。ビルのゆれのような、ゆっくりした（\qty{1}{Hz}くらい）ゆれでは変位を、機械の軸のような数十 Hz〜数百 Hz のゆれでは速度を、歯車や軸受けの傷のような数 kHz 以上のゆれでは加速度をはかるのが、ふつうです。機械の振動の良い・悪いの目安は、速度（mm/s）で決められていることが多く、規格に「良い」「注意」「危ない」の区切りの値が示されています。

%% ---------------------------------------------------------------
\bksubsection{7}{ゆれをはかるセンサーのしくみ}

\leavevmode\label{ph:波-10:sec7}ゆれをはかりたい物（機械や地面）にセンサーをとりつけると、センサーの入れ物（ケース）も、いっしょにゆれます。ケースの中に、ばねとダンパーでつるしたおもりを入れておきます（図\ref{ph:波-10:fig-sensor}）。ケースがゆれると、おもりは少しおくれてついていくので、ケースとおもりの間がずれます。このずれ（ばねの伸び縮み）を電気の信号にかえて、ゆれをはかるのです。

\begin{bkfigure}
\begin{tikzpicture}[x=1mm, y=1mm]
  % ケース
  \draw[bkInk, line width=1.1pt] (0,0) rectangle (44,42);
  \fill[bkGray!20] (0,-4) rectangle (44,0);
  \draw[bkInk, line width=0.8pt] (0,-4) rectangle (44,0);
  \node[rxk lab, anchor=north] at (22,-5) {はかりたい物（機械・地面）};
  \node[rxk lab, anchor=south west] at (1,42.5) {ケース};
  % おもり
  \filldraw[fill=bkBlue!30, draw=bkInk, line width=0.8pt] (8,20) rectangle (36,30);
  \node[rxk lab] at (22,25) {おもり $m$};
  \rxkspring{(14,0)}{20}
  \rxkdamper{(30,0)}{20}
  \node[rxk lab, anchor=east] at (11,10) {ばね};
  \node[rxk lab, anchor=west] at (33,8) {ダンパー};
  % 変位
  \draw[rxk arr, bkRed!75!black] (52,-2) -- (52,8);
  \node[rxk lab, anchor=west, text=bkRed!70!black] at (54,3) {ケースのゆれ $y$};
  \draw[rxk arr, bkOrange!85!black] (52,25) -- (52,35);
  \node[rxk lab, anchor=west, text=bkOrange!70!black, align=left] at (54,28) {おもりの、ケースに\\対するずれ $z$};
\end{tikzpicture}
\bkfigcaption{振動センサーの模型。ケースが $y$ だけゆれると、おもりはケースに対して $z$ だけずれる。$z$ を電気の信号にしてはかる。}
\label{ph:波-10:fig-sensor}
\end{bkfigure}

\noindent\textbf{運動方程式}\par\nobreak
おもりの、地面から見た位置は $y+z$ です。おもりにはたらく力は、ばねの力（$z$ に比例）とダンパーの力（ずれる速さ $\dot z$ に比例）だけです。1 kg あたりで書くと、波-1 ⑦と同じく、固有角振動数 $\omega_{\mathrm n}$ と減衰比 $\zeta$ を使って
\[ \ddot y+\ddot z=-\omega_{\mathrm n}^{\,2}z-2\zeta\omega_{\mathrm n}\dot z \]
（左辺はおもりの加速度）。ケースの加速度 $\ddot y$ を右辺へ移すと
\pheq{\ddot z+2\zeta\omega_{\mathrm n}\dot z+\omega_{\mathrm n}^{\,2}z=-\ddot y}{m/s$^2$}
\label{ph:波-10:sensoreq}です。これは、波-2 の「ゆすり続けるばねとおもり」の式で、ゆする力の代わりに、ケースの加速度 $-\ddot y$ が入った形です。

\noindent\textbf{ずれの大きさ}\par\nobreak
ケースが、角振動数 $\omega$、加速度の振幅 $a_0$ でゆれ続けるとします。波-2 ④（6-13⑥）と同じ計算で、ずれの振幅は
\pheq{|z|=\frac{a_0}{\omega_{\mathrm n}^{\,2}}\,H,\qquad H=\frac{1}{\sqrt{\bigl(1-(\omega/\omega_{\mathrm n})^2\bigr)^2+\bigl(2\zeta\,\omega/\omega_{\mathrm n}\bigr)^2}}}{m ＝ m/s$^2$ ÷ (1/s$^2$)}
\label{ph:波-10:H}になります。$H$ は「ゆっくり押したときに比べて、ずれが何倍になるか」を表す\textbf{ゆれの倍率}で、振動数の比 $\omega/\omega_{\mathrm n}$ と減衰比 $\zeta$ だけで決まります。

\noindent\textbf{加速度計：$\omega\ll\omega_{\mathrm n}$ で使う}\par\nobreak
ゆれの振動数が、センサーの固有振動数よりずっと低いとき（$\omega/\omega_{\mathrm n}\to0$）、$H\to1$ なので
\[ |z|\approx\frac{a_0}{\omega_{\mathrm n}^{\,2}} \]
で、ずれはケースの\textbf{加速度}に比例します。これが\textbf{加速度計}です。\bkidx{かそくどけい}{加速度計}ばねを硬く（$\omega_{\mathrm n}$ を大きく）して、はかりたいゆれよりずっと高い固有振動数にしておけば、ずれから加速度がわかります。

\noindent\textbf{変位計（地震計）：$\omega\gg\omega_{\mathrm n}$ で使う}\par\nobreak
反対に、ばねをとてもやわらかくして、固有振動数をはかりたいゆれよりずっと低くします。ケースの変位の振幅を $Y$ とすると $a_0=\omega^2Y$ なので
\[ |z|=\Bigl(\frac{\omega}{\omega_{\mathrm n}}\Bigr)^2H\,Y\ \to\ Y\qquad(\omega/\omega_{\mathrm n}\to\infty) \]
です。おもりは、速いゆれについていけず、ほとんどその場に止まっています。だから、ケースとおもりのずれが、ケースの\textbf{変位}そのものになります。昔の地震計は、重いおもりを長い周期のふりこでつるして、地面のゆれを紙に書かせていました。これが\textbf{変位計}（地震計）の考え方です。\bkidx{じしんけい}{地震計}

\Needspace{14\baselineskip}
\begin{phflow}
\phflowrow{$\ddot z+2\zeta\omega_{\mathrm n}\dot z+\omega_{\mathrm n}^{\,2}z=-\ddot y$}{ゆれ続ける\\（6-13 の特解）}{$|z|=\dfrac{a_0}{\omega_{\mathrm n}^{\,2}}H$}
\phflowrow{$|z|=\dfrac{a_0}{\omega_{\mathrm n}^{\,2}}H$}{$\omega\ll\omega_{\mathrm n}$\\（$H\approx1$）}{$|z|\approx\dfrac{a_0}{\omega_{\mathrm n}^{\,2}}$\\加速度計}
\phflowrow{$|z|=\Bigl(\dfrac{\omega}{\omega_{\mathrm n}}\Bigr)^2HY$}{$\omega\gg\omega_{\mathrm n}$}{$|z|\approx Y$\\変位計（地震計）}
\end{phflow}

\begin{bkfigure}
\begin{tikzpicture}
\begin{axis}[bkplot, xmode=log, width=108mm, height=56mm, xmin=0.05, xmax=1.6, ymin=0, ymax=2.6, xtick={0.05,0.1,0.2,0.5,1}, xticklabels={0.05,0.1,0.2,0.5,1}, ytick={0.5,1,1.5,2,2.5}, xlabel={$\omega/\omega_{\mathrm n}$}, ylabel={$H$}, axis x line*=bottom, axis y line*=left]
  \fill[bkGreen!12] (axis cs:0.05,0.95) rectangle (axis cs:1.6,1.05);
  \addplot[bkRed!80!black, line width=1.4pt, densely dotted, domain=0.05:1.6, samples=200] {1/sqrt((1-x^2)^2+(0.2*x)^2)};
  \addplot[bkOrange, line width=1.5pt, dashed, domain=0.05:1.6, samples=150] {1/sqrt((1-x^2)^2+(0.8*x)^2)};
  \addplot[bkcurve, domain=0.05:1.6, samples=150] {1/sqrt((1-x^2)^2+(1.4*x)^2)};
  \node[rxk lab, anchor=west, text=bkRed!70!black] at (axis cs:0.06,2.3) {点線：$\zeta=0.1$（$\omega=\omega_{\mathrm n}$ で5倍）};
  \node[rxk lab, anchor=west, text=bkOrange!80!black] at (axis cs:0.06,1.95) {破線：$\zeta=0.4$};
  \node[rxk lab, anchor=west, text=bkBlue!80!black] at (axis cs:0.06,1.6) {実線：$\zeta=0.7$（平らな所が広い）};
  \node[rxk lab, anchor=north west, text=bkGreen!40!black] at (axis cs:0.06,0.9) {色の帯：ずれが $\pm5$\,\% 以内};
\end{axis}
\end{tikzpicture}
\bkfigcaption{加速度計のゆれの倍率 $H$。$H=1$ からのずれが、感度のくるい。$\zeta=0.7$ くらいだと、$\omega/\omega_{\mathrm n}\approx0.6$ まで、くるいが5\,\% 以内におさまる。}
\label{ph:波-10:fig-H}
\end{bkfigure}

\noindent\textbf{なぜ $\zeta\approx0.7$ にするのか}\par\nobreak
図\ref{ph:波-10:fig-H}のように、ダンパーが弱い（$\zeta=0.1$）と、固有振動数の近くで共振して、$H$ が大きくなりすぎます。ダンパーが強すぎると、$H$ は1より早く小さくなります。ちょうどよいのが $\zeta=1/\sqrt2\approx0.7$ です。このとき $(2\zeta)^2=2$ なので、$H$ の分母の中は
\[ \bigl(1-(\omega/\omega_{\mathrm n})^2\bigr)^2+2(\omega/\omega_{\mathrm n})^2=1+(\omega/\omega_{\mathrm n})^4 \]
となり、$(\omega/\omega_{\mathrm n})^2$ の項が消えます。$\omega/\omega_{\mathrm n}=0.5$ でも $(\omega/\omega_{\mathrm n})^4=0.0625$ と小さいので、$H\approx0.97$ です（\bkreidairef{11}）。広い振動数の範囲で、感度が平らになるのです。ただし、工場でよく使う圧電式の加速度計は、中にダンパーを入れにくく、$\zeta$ がとても小さいセンサーです。そこで、固有振動数の約3分の1までの振動数で使います（$\omega/\omega_{\mathrm n}=1/3$ で、$\zeta\approx0$ なら $H=1/(1-1/9)\approx1.13$。くるいは約12\,\%）。

\begin{phcode}[Python]
import math
def H(r, z):                       \# ゆれの倍率
    return 1 / math.sqrt((1 - r*r)**2 + (2*z*r)**2)
print('w/wn ', '  '.join(f'z={z}' for z in (0.1, 0.4, 0.7)))
for r in (0.1, 0.2, 0.3, 0.5, 0.7, 1.0):
    s = [f'{H(r, z):5.2f}' for z in (0.1, 0.4, 0.7)]
    print(f'{r:4.1f} ', '  '.join(s))
\end{phcode}
\begin{phcodeout}
w/wn  z=0.1  z=0.4  z=0.7
 0.1   1.01   1.01   1.00
 0.2   1.04   1.03   1.00
 0.3   1.10   1.06   1.00
 0.5   1.32   1.18   0.97
 0.7   1.89   1.32   0.91
 1.0   5.00   1.25   0.71
\end{phcodeout}
プログラムの \texttt{r} は $\omega/\omega_{\mathrm n}$、\texttt{z} は $\zeta$ です。$\zeta=0.7$ の列は、0.3 まで1.00 のまま、0.5 でも0.97 です。$\zeta=0.1$ の列は、0.3 ですでに1割大きくなっています。

\begin{phmijika}[スマホの中の加速度センサー]
スマホの加速度センサーは、髪の毛の太さくらいの小さなおもりと、くし形の電極でできていて、半導体の加工の技術（MEMS）で作られています。固有振動数は、人の動きや車のゆれ（数十 Hz まで）よりずっと高く作られているので、それらのゆれは、$\omega\ll\omega_{\mathrm n}$ の「加速度計の領分」に入ります。止まっていても、重力の分だけおもりが下にずれるので、どちらが下かもわかります。
\end{phmijika}

%% ---------------------------------------------------------------
\bksubsection{8}{回る機械のゆれで、故障を見つける}

\leavevmode\label{ph:波-10:sec8}工場のポンプやモーター、発電所のタービンには、振動センサーがとりつけられています。とったゆれのデータを FFT（9-4⑬）でスペクトルに分けると、どの振動数の成分が大きいかがわかります。9-4・9-12④で読んだように、回転数の何倍の成分が大きいかで、悪い所の見当がつきます。

1分に\num{1800}回転する軸は、1秒に $1800\div60=30$ 回転なので、回転の周波数は $f_{\mathrm r}=\qty{30}{Hz}$ です。
\begin{itemize}[leftmargin=1.2\zw, itemsep=2pt, topsep=2pt, label={\color{bkTheme}●}]
\item \textbf{1倍（\qty{30}{Hz}）が大きい}：回る部分の重さのかたより（不つり合い）。1回転に1回、重い所が外へ引っぱるので、回転と同じ振動数でゆれる（波-2）。
\item \textbf{2倍（\qty{60}{Hz}）が大きい}：モーターとポンプの軸のつなぎ目のずれ（芯ずれ）。1回転に2回、押し合うことが多い。
\item \textbf{回転数の整数倍でない線}：軸受けの傷。
\end{itemize}

\noindent\textbf{軸受けの傷の周波数}\par\nobreak
玉軸受けは、外輪と内輪の間に、$n$ 個の玉がはさまって転がる部品です（図\ref{ph:波-10:fig-jikuuke}）。外輪は止まっていて、内輪が軸といっしょに回ります。外輪に傷があると、玉が傷の上を通るたびに「コツン」とたたきます。玉の中心を結ぶ円の直径を $D_{\mathrm p}$、玉の直径を $D_{\mathrm b}$ とすると、玉の列（保持器）は、軸の回転の $\tfrac12(1-D_{\mathrm b}/D_{\mathrm p})$ 倍の速さで回ります（下のせいかくに言うと）。外輪の傷を1秒にたたく回数は、玉の数×（保持器の回転数）なので
\pheq{f_{\text{外輪}}=\frac{n}{2}\,f_{\mathrm r}\Bigl(1-\frac{D_{\mathrm b}}{D_{\mathrm p}}\Bigr)}{Hz ＝ Hz}
です。玉9個・$D_{\mathrm b}/D_{\mathrm p}=0.2$・$f_{\mathrm r}=\qty{30}{Hz}$ なら $4.5\times30\times0.8=\qty{108}{Hz}$ で（\bkreidairef{14}）、30 の整数倍ではありません。スペクトルに\qty{108}{Hz}とその整数倍の線が出てきたら、外輪の傷をうたがいます。内輪の傷は、内輪が玉の列を追いこしながら回るので、$\frac{n}{2}f_{\mathrm r}(1+D_{\mathrm b}/D_{\mathrm p})$ になります。

\begin{bkfigure}
\begin{tikzpicture}[x=1mm, y=1mm]
  % 軸受け
  \begin{scope}[shift={(22,20)}]
  \filldraw[fill=bkGray!30, draw=bkInk, line width=0.8pt] (0,0) circle (19);
  \filldraw[fill=white, draw=bkInk, line width=0.8pt] (0,0) circle (15.5);
  \filldraw[fill=bkGray!30, draw=bkInk, line width=0.8pt] (0,0) circle (8.5);
  \filldraw[fill=white, draw=bkInk, line width=0.8pt] (0,0) circle (5);
  \foreach \k in {0,...,8} {\filldraw[fill=bkBlue!35, draw=bkInk, line width=0.6pt] ({90+\k*40}:12) circle (3.3);}
  \draw[bkRed!75!black, line width=2pt] (-60:15.4) -- (-60:17);
  \draw[rxk dim] (0,0) -- (0:12);
  \draw[rxk arr] (120:6.5) arc (120:60:6.5);
  \end{scope}
  \node[rxk lab, anchor=west] at (44,36) {外輪（止まっている）};
  \node[rxk lab, anchor=west] at (44,30) {玉（$n=9$、直径 $D_{\mathrm b}$）};
  \node[rxk lab, anchor=west] at (44,24) {内輪（軸と回る）};
  \node[rxk lab, anchor=west] at (44,18) {太い短い線：傷};
  \node[rxk lab, anchor=west] at (44,12) {両矢印：$D_{\mathrm p}/2$};
  % スペクトル
  \begin{scope}[shift={(92,0)}]
  \draw[bkInk, line width=0.7pt, -{Stealth[length=3pt]}] (0,0) -- (38,0);
  \draw[bkInk, line width=0.7pt, -{Stealth[length=3pt]}] (0,0) -- (0,34);
  \node[rxk lab, anchor=south] at (0,34.5) {大きさ};
  \foreach \f/\h/\t in {30/24/30,60/10/60,108/18/108,216/8/216} {\draw[bkBlue!80!black, line width=1.6pt] ({\f*0.15},0) -- ({\f*0.15},\h); \node[rxk lab, anchor=north] at ({\f*0.15},-0.8) {\t};}
  \node[rxk lab, anchor=north] at (19,-5) {振動数（Hz）};
  \end{scope}
\end{tikzpicture}
\bkfigcaption{左：玉軸受け（正面から見た図）。右：外輪に傷があるポンプ（毎分1800回転）のスペクトルの例。30・60 Hz は回転の1倍・2倍、108・216 Hz は外輪の傷。}
\label{ph:波-10:fig-jikuuke}
\end{bkfigure}

\noindent\textbf{せいかくに言うと}\par\nobreak
外輪が止まり、内輪が速さ $V$ で回っているとき、玉は「外輪にふれる点が止まり、内輪にふれる点が $V$ で動く」転がり方をします。玉の中心は、その間のまん中なので、速さは $V/2$ です。内輪の表面の速さは $V=\pi(D_{\mathrm p}-D_{\mathrm b})f_{\mathrm r}$ なので、玉の中心が円（直径 $D_{\mathrm p}$）を回る回数は $\tfrac12(1-D_{\mathrm b}/D_{\mathrm p})f_{\mathrm r}$ です。玉が斜めに押されて転がる軸受け（接触角があるもの）では、$D_{\mathrm b}/D_{\mathrm p}$ に接触角の cos がかかります。本当の機械では、玉が少しすべるので、計算より数\,\% ずれることがあります。

%% ---------------------------------------------------------------
\bksubsection{9}{たたいて調べる：インパルスハンマーと打音検査}

\leavevmode\label{ph:波-10:sec9}スイカをたたいて熟れ具合を確かめるように、物をたたいて、その響き方で中の様子を知る方法があります。

\noindent\textbf{インパルスハンマー}\par\nobreak
機械の部品や建物のはりを、力をはかれるハンマー（\textbf{インパルスハンマー}）で軽くたたき、ゆれを加速度計ではかる試験があります。\bkidx{いんぱるすはんまー}{インパルスハンマー}9-6⑦で学んだように、ごく短い時間だけはたらく力は、デルタ関数 $\delta(t)$ で表せます。そのときのゆれ（\textbf{インパルス応答}）には、その物の固有振動数の成分が全部ふくまれています（図\ref{ph:波-10:fig-hammer}）。インパルス応答がわかれば、9-6⑧のたたみこみで、どんな力を加えたときのゆれも計算できます。インパルス応答をスペクトルに分けると、山の位置で固有振動数が、山の鋭さで減衰比がわかります（波-2 ⑤）。

\begin{bkfigure}
\begin{tikzpicture}
\begin{axis}[bkplot, width=108mm, height=46mm, xmin=-0.02, xmax=0.42, ymin=-1.15, ymax=1.65, xtick={0.1,0.2,0.3,0.4}, xticklabels={0.1,0.2,0.3,0.4}, ytick={-1,0,1}, xlabel={$t$（s）}, axis x line*=bottom, axis y line*=left]
  \draw[bkGray, line width=0.5pt] (axis cs:0,0) -- (axis cs:0.41,0);
  \addplot[bkcurve, domain=0:0.4, samples=400] {exp(-2*pi*25*0.02*x)*sin(deg(2*pi*25*x))};
  \draw[bkRed!80!black, line width=1.6pt] (axis cs:0.003,0) -- (axis cs:0.003,1.35);
  \node[rxk lab, anchor=west, text=bkRed!70!black] at (axis cs:0.006,1.4) {太い縦線：たたいた力（ごく短い）};
  \node[rxk lab, anchor=south west, text=bkBlue!80!black] at (axis cs:0.16,0.75) {細い線：ゆれ（\qty{25}{Hz}で、だんだん小さくなる）};
\end{axis}
\end{tikzpicture}
\bkfigcaption{はりを、インパルスハンマーでたたいたときのゆれ（インパルス応答）の例。固有振動数\qty{25}{Hz}、減衰比0.02 として計算。山と山の間は\qty{40}{ms}。}
\label{ph:波-10:fig-hammer}
\end{bkfigure}

\noindent\textbf{打音検査}\par\nobreak
トンネルの壁や橋のコンクリートは、点検のとき、人がハンマーでたたいて、音を聞いて調べます。これを\textbf{打音検査}といいます。\bkidx{だおんけんさ}{打音検査}中が詰まった健全な所は「コンコン」と高く澄んだ音がしますが、中にすき間ができて表面が浮いている所は「ポコポコ」と低く濁った音がします。浮いた部分は、周りから切り離された薄い板のようになり、やわらかく（固有振動数が低く）なるからです（波-6）。浮いた部分は、はがれて落ちる危険があるので、たたき落とすか、補修します。最近は、たたいた音をマイクで録音し、コンピューターでスペクトルを調べて判定する研究も進んでいます。

%% ---------------------------------------------------------------
\bksubsection{10}{ひずみゲージで振動をはかる}

\leavevmode\label{ph:波-10:sec10}力-9 ⑩では、部品にはりつけた\textbf{ひずみゲージ}の電気の抵抗の変化から、ひずみと応力を測りました。\bkidx{ひずみげーじ}{ひずみゲージ}ひずみゲージは、ゆっくりした力だけでなく、振動による速いひずみの変化もはかれます。抵抗の変わる割合は、ひずみ $\varepsilon$ のおよそ2倍（ゲージ率 約2）です。
\[ \frac{\Delta R}{R}\approx2\,\varepsilon \]

\begin{bkfigure}
\begin{tikzpicture}[x=1mm, y=1mm]
  \fill[bkGray!35] (0,-4) rectangle (6,16);
  \draw[bkInk, line width=0.8pt] (6,-4) -- (6,16);
  \foreach \yy in {-3,0,...,15} {\draw[bkInk, line width=0.5pt] (0,\yy) -- (-2,{\yy-2});}
  \filldraw[fill=bkBlue!25, draw=bkInk, line width=0.8pt] (6,4) -- (86,4) -- (86,8) -- (6,8) -- cycle;
  \draw[bkBlue!70!black, line width=0.8pt, dashed] (6,8) .. controls (40,8) and (60,6) .. (86,1.5);
  \draw[bkBlue!70!black, line width=0.8pt, dashed] (6,8) .. controls (40,8) and (60,10) .. (86,14.5);
  \filldraw[fill=bkOrange!50, draw=bkInk, line width=0.6pt] (10,8) rectangle (20,9.4);
  \node[rxk lab, anchor=south west] at (19,9.6) {ひずみゲージ};
  \draw[rxk arr] (95,3) -- (95,14);
  \draw[rxk arr] (95,14) -- (95,3);
  \node[rxk lab, anchor=west] at (97,8.5) {先がゆれる};
  \node[rxk lab, anchor=north] at (46,2.5) {片持ちのはり（破線はゆれたときの形）};
\end{tikzpicture}
\bkfigcaption{片持ちのはりの根もとにひずみゲージをはる。はりが上下にゆれると、根もとの上の面が、伸びたり縮んだりを繰り返す。}
\label{ph:波-10:fig-gauge}
\end{bkfigure}

片持ちのはりがゆれるとき、根もとの表面は、伸びと縮みを繰り返します（力-10 で学んだように、曲げのひずみは根もとで最も大きい）。\qty{120}{\ohm}のゲージで、ひずみの振幅が $200\times10^{-6}$（200 マイクロストレイン）なら、抵抗の変化の振幅は $120\times2\times200\times10^{-6}=\qty{0.048}{\ohm}$ です（\bkreidairef{15}）。この小さな変化を、ブリッジ回路（電-5）で電圧の信号にして、記録します。ひずみの記録の山の数を数えれば振動数が、大きさから応力（$E\varepsilon$）がわかるので、橋や車の部品が、ゆれのたびにどれだけの応力を受けているかを調べられます（疲労の点検。力-8）。

\begin{phmada}[コウモリは、どうして仲間のこだまとまちがえない？]
洞くつから、何千匹ものコウモリがいっしょに飛び出すとき、周りは、仲間が出す超音波と、そのこだまでいっぱいです。それでも、コウモリはぶつからずに飛び、自分のこだまで虫をつかまえています。

自分の声の高さを少しずらして、仲間の声と重ならないようにする種があると報告されています。一方で、声の高さをずらさないという研究もあり、意見は分かれています。そして、たくさんの音がまざった中から、どうやって自分のこだまだけを選び出しているのか、そのしくみの全体はまだわかっていません。耳と脳がどのように音を処理しているのか、小さなマイクを背負わせたコウモリの記録などを使って、研究が続いています。この研究は、たくさんの車がいっしょに使うレーダーやセンサーが、互いにじゃまし合わないようにする技術のヒントにもなると期待されています。
\end{phmada}

%% ---------------- たいせつ・ことば ----------------
\begin{taisetsu}
\begin{itemize}[leftmargin=1.2\zw, itemsep=2pt, topsep=0pt, label={\color{bkNavy}●}]
\item こだまの距離 $L=ct/2$（往復）。空気\qty{343}{m/s}、海水 約\qty{1500}{m/s}、鋼 約\qty{5900}{m/s}。
\item 振動数を上げると、細かく見分けられる（波長が短い）が、早く弱まって深くまで届かない。
\item 傷や空気のすき間は、$Z$ が大きく違うので、ほとんど全部はね返す（$r\approx-1$）。
\item $v_{\max}=2\pi fA$、$a_{\max}=(2\pi f)^2A$。振動数2倍で、速度2倍、加速度4倍。
\item 振動センサー：$|z|=(a_0/\omega_{\mathrm n}^{\,2})H$。$\omega\ll\omega_{\mathrm n}$ で加速度計、$\omega\gg\omega_{\mathrm n}$ で変位計。$\zeta\approx0.7$ で感度が平ら。
\item 回転の1倍は不つり合い、2倍は芯ずれ、整数倍でない線は軸受けの傷。
\end{itemize}
\end{taisetsu}

\begin{kotoba}\small
\bkword{エコーロケーション}{自分で出した音のこだまで、周りの物の場所を知ること。}
\bkword{超音波探傷・パルス反射法}{短い超音波を入れて、傷や底面からのこだまの時間で中を調べる方法。}
\bkword{圧電効果}{押すと電気が出て、電気をかけるとのび縮みする性質。}
\bkword{加速度計・変位計}{ばねとおもりのずれで、加速度（$\omega\ll\omega_{\mathrm n}$）や変位（$\omega\gg\omega_{\mathrm n}$）をはかるセンサー。}
\bkword{インパルス応答}{ごく短い力でたたいたときのゆれ。固有振動数と減衰がわかる。}
\bkword{打音検査}{たたいた音の高さや濁りで、中のすき間を見つける点検。}
\end{kotoba}
\bkkigocards{\bkkigocard{$d_{33}$}{ディー・さん・さん}{圧電定数〔C/N〕。押す力1 N あたりにたまる電荷。33 は押す向きと電気の出る向きの番号}{piezoelectric constant}{波-10}}{\bkkigocard{$H$}{エイチ}{ゆれの倍率。ゆっくり押したときに比べて、ずれが何倍か。$\omega/\omega_{\mathrm n}$ と $\zeta$ で決まる。波-11 の $H$ は情報のエントロピー}{magnification factor}{波-10}}

\begin{senmonsho}[振動計測の本の書き方]
\begin{bkhonbun}
相対変位形の振動ピックアップの振幅比は $|z/y|=\eta^2/\sqrt{(1-\eta^2)^2+(2\zeta\eta)^2}$ で表される。$\eta\ll1$ では $z\approx-\ddot y/\omega_n^{\,2}$ となり加速度計として、$\eta\gg1$ では $z\approx-y$ となり変位計として働く。減衰比を $\zeta\approx0.7$ にとると、使用できる周波数範囲が最も広くなる。
\end{bkhonbun}
\bkyaku{$\eta$（イータ）は振動数の比 $\omega/\omega_{\mathrm n}$。この本では文字を作らず $\omega/\omega_{\mathrm n}$ と書いた。振動ピックアップは振動センサーのこと。$|z/y|$ は⑦の $(\omega/\omega_{\mathrm n})^2H$。$-$ がつくのは、ケースとおもりのずれの向きが、ケースの加速度と反対だから。}
\end{senmonsho}

%% ---------------- 例題 ----------------
\Needspace{14\baselineskip}
\bksection{reidai}{例題でたしかめよう}

\noindent 空気の音の速さは\qty{343}{m/s}、海水は\qty{1500}{m/s}、鋼の縦波は\qty{5900}{m/s}、体の軟らかい所は\qty{1540}{m/s}とします。

\begin{phreidai}[コウモリと虫の距離]
\label{ph:波-10:ex1}コウモリが超音波を出してから、虫からのこだまが返ってくるまで6.0 ms でした。虫までの距離を求めましょう。コウモリが\qty{100}{kHz}の超音波を使うとき、その波長はいくらですか。
\phbamen
\bkhatenaref の場面です。コウモリは、こだまの時間で距離を知る。
\phmushi
コウモリと虫が動く分（6 ms の間に、コウモリは数 cm しか進まない）。風。気温による音の速さの違い（\qty{20}{\celsius}とする）。
\phzu
虫までの距離 $L$、音の速さ $c=\qty{343}{m/s}$、往復の時間 $t=\qty{6.0e-3}{s}$、振動数 $f$、波長 $\lambda$。
\phriyu
音は虫まで行って帰るので、道のりは $2L$（①）。波長は $c=f\lambda$（波-5 ②）から。
\phtateru
\[ L=\frac{c\,t}{2},\qquad \lambda=\frac{c}{f} \]
\phtoku
\[ L=\frac{343\times6.0\times10^{-3}}{2}=\frac{2.058}{2}\approx\qty{1.0}{m} \]
\[ \lambda=\frac{343}{100\times10^3}=\qty{3.43e-3}{m}\approx\qty{3.4}{mm} \]
\phtashikame
単位：m/s × s ＝ m。大きさ：コウモリが虫をつかまえる直前の距離として自然。波長\qty{3.4}{mm}は、小さな虫くらい（波長くらいの物まで見分けられる：②）。極端な場合：往復を忘れると\qty{2.1}{m}になり、2倍にずれる。
\kotae 約\qty{1.0}{m}、波長は約\qty{3.4}{mm}
\end{phreidai}
\begin{jibun}
コウモリのこだまが2.0 ms 後に返ってきました。虫までの距離は何 m ですか。また、\qty{40}{kHz}の超音波の波長は何 mm ですか。
\bkjans{$343\times0.0020\div2\approx\qty{0.34}{m}$。$343\div40000\approx\qty{0.0086}{m}$（約\qty{8.6}{mm}）}
\end{jibun}

\begin{reidai}[車の障害物センサー]
\label{ph:波-10:ex2}車の後ろのセンサーが出した超音波のこだまが、2.9 ms 後に返ってきました。壁までの距離を求めましょう。
\par\bkstep{理想化} 音はまっすぐ壁へ行き、まっすぐ返る。気温は\qty{20}{\celsius}。
\par\kangaekata 往復なので $L=ct/2$（①）。
\par\shiki $L=343\times0.0029\div2=0.99470\div2$
\par\bkstep{確かめ} 約\qty{0.5}{m}は、「ピー」と鳴り続けて止まる目安くらいの距離。
\par\kotae 約\qty{0.50}{m}
\end{reidai}
\begin{jibun}
壁まで\qty{0.30}{m}のとき、こだまは何 ms 後に返ってきますか。
\bkjans{$2\times0.30\div343\approx\qty{1.7e-3}{s}$（約1.7 ms）}
\end{jibun}

\begin{phreidai}[おなかの中を見る超音波の振動数]
\label{ph:波-10:ex3}おなかの中を、\qty{5}{MHz}の超音波で調べます。体の中の波長を求めましょう。体の中の減衰を 0.5\,dB/(cm·MHz) とすると、深さ\qty{10}{cm}の所からのこだまは、行って帰るまでに何 dB 弱まりますか。
\phbamen
超音波の検査で、見分けられる大きさと、届く深さを見積もる。
\phmushi
音が広がって弱まる分、はね返る割合（減衰だけを考える）。場所による音の速さの違い。
\phzu
振動数 $f=\qty{5}{MHz}$、音の速さ $c=\qty{1540}{m/s}$、波長 $\lambda$、深さ\qty{10}{cm}（往復\qty{20}{cm}）。
\phriyu
波長は $\lambda=c/f$。減衰は「1 cm・1 MHz あたり 0.5 dB」なので、（0.5）×（MHz の数）×（往復の cm）で dB になる（②）。dB はたし算できる（3-12⑦）。
\phtateru
\[ \lambda=\frac{c}{f},\qquad (\text{弱まり})=0.5\times5\times20 \]
\phtoku
\[ \lambda=\frac{1540}{5\times10^6}=\qty{3.08e-4}{m}\approx\qty{0.31}{mm},\qquad 0.5\times5\times20=\qty{50}{dB} \]
\phtashikame
大きさ：50\,dB は、強さで10万分の1。それでもこだまは聞きとれるが、もっと深く見るには振動数を下げる。\qty{10}{MHz}にすると、波長は半分になるが、弱まりは2倍の100\,dB になる。
\kotae 波長は約\qty{0.31}{mm}、弱まりは約50\,dB
\end{phreidai}
\begin{jibun}
\qty{3.5}{MHz}の超音波で、深さ\qty{8}{cm}の所を見るとき、体の中の波長と、往復の弱まりを求めましょう。
\bkjans{$1540\div(3.5\times10^6)\approx\qty{0.44}{mm}$。$0.5\times3.5\times16=\qty{28}{dB}$}
\end{jibun}

\begin{phreidai}[鋼の板の中の傷の深さ]
\label{ph:波-10:ex4}厚さ\qty{100}{mm}の鋼の板を、パルス反射法で調べました。画面に、出した音から\qty{20}{\micro s}後にこだまが出ました。傷の深さを求めましょう。また、底面のこだまは何 μs 後に出るはずですか。
\phbamen
橋の鋼の板の中に、割れがないかを調べる。
\phmushi
探触子と板の間の油の層の厚さ（とても薄い）。板の中の場所による音の速さの違い。
\phzu
傷の深さ $L$、鋼の縦波の速さ $c=\qty{5900}{m/s}$、こだまの時間 $t=\qty{20e-6}{s}$、板の厚さ $\qty{0.100}{m}$。
\phriyu
傷まで行って帰るので、$L=ct/2$（③）。底面までの時間は、同じ式を時間について解いた $t=2L/c$。
\phtateru
\[ L=\frac{c\,t}{2},\qquad t_{\text{底}}=\frac{2\times0.100}{c} \]
\phtoku
\[ L=\frac{5900\times20\times10^{-6}}{2}=\qty{0.0592}{m}\approx\qty{59}{mm} \]
\[ t_{\text{底}}=\frac{0.200}{5900}\approx\qty{3.4e-5}{s}=\qty{34}{\micro s} \]
\phtashikame
単位：m/s × s ＝ m。大きさ：傷の深さ\qty{59}{mm}は、板の厚さ\qty{100}{mm}より浅いので、こだまは底面（\qty{34}{\micro s}）より先に出る。順番も合う。
\kotae 傷の深さは約\qty{59}{mm}、底面のこだまは約\qty{34}{\micro s}後
\end{phreidai}
\begin{jibun}
鋼の板で、傷のこだまが\qty{25}{\micro s}後に出ました。傷の深さは何 mm ですか。板の厚さが\qty{80}{mm}なら、底面のこだまは何 μs 後ですか。
\bkjans{$5900\times25\times10^{-6}\div2\approx\qty{0.074}{m}$（\qty{74}{mm}）。底面は $2\times0.080\div5900\approx\qty{27}{\micro s}$（傷の\qty{25}{\micro s}より後）}
\end{jibun}

\begin{phreidai}[水から鋼に入る超音波]
\label{ph:波-10:ex5}水の中に沈めた鋼の板に、超音波を垂直に当てます。水の $Z$ は\qty{1.48e6}{Pa.s/m}、鋼の $Z$ は\qty{4.63e7}{Pa.s/m}です。表面での反射係数と、はね返る強さ・鋼に入る強さの割合を求めましょう。鋼の中の傷（空気のすき間。$Z$ は\qty{413}{Pa.s/m}）では、どうですか。
\phbamen
水浸法の探傷で、超音波がどれだけ鋼に入るかを見積もる。
\phmushi
斜めに当たる分（垂直とする）。水や鋼の中で弱まる分。傷の大きさ（傷は超音波のはばより大きいとする）。
\phzu
入る側の $Z_1$、向こう側の $Z_2$、音圧の反射係数 $r$、はね返る強さの割合 $r^2$、入る強さの割合 $1-r^2$。
\phriyu
境目で音圧と動く速さがつながる、という条件から導いた反射係数（波-7 ⑤）を、音圧で数えた形（波-7 ④の $r_{\text{押}}$）で使う。強さは音圧の2乗に比例するので、強さの割合は $r^2$。エネルギーは保存するので、残りが入る。
\phtateru
\[ r=\frac{Z_2-Z_1}{Z_2+Z_1},\qquad (\text{はね返る})=r^2,\qquad (\text{入る})=1-r^2 \]
\phtoku
水→鋼：
\[ r=\frac{4.63\times10^7-1.48\times10^6}{4.63\times10^7+1.48\times10^6}=\frac{4.48\times10^7}{4.78\times10^7}\approx0.94 \]
はね返る強さ $0.94^2\approx0.88$、入る強さ $1-0.88=0.12$。

鋼→空気（傷）：
\[ r=\frac{413-4.63\times10^7}{413+4.63\times10^7}\approx-1.00 \]
はね返る強さ $r^2\approx1.00$。ほぼ全部はね返る。
\phtashikame
単位：Pa·s/m ÷ Pa·s/m で単位なし。大きさ：$Z$ の差が大きいほど $|r|$ は1に近い。極端な場合：$Z_1=Z_2$ なら $r=0$ で、全部入る（ゼリーのはたらき：波-7）。
\kotae 水→鋼：$r\approx0.94$、はね返る88\,\%、入る12\,\%。傷（空気）：$r\approx-1$ で、ほぼ全部はね返る
\end{phreidai}
\begin{jibun}
体（$Z=\qty{1.63e6}{Pa.s/m}$）から空気（\qty{413}{Pa.s/m}）へ向かう超音波の反射係数を求めましょう。ゼリー（\qty{1.5e6}{Pa.s/m}）から体へ入るときはどうですか。
\bkjans{体→空気 $r=(413-1.63\times10^6)$ $\div(413+1.63\times10^6)$ $\approx-1.00$。ゼリー→体 $r=(1.63-1.5)\div(1.63+1.5)$ $\approx0.04$（強さで0.2\,\% しかはね返らない）}
\end{jibun}

\begin{reidai}[配管の厚さをはかる]
\label{ph:波-10:ex6}鋼の配管に厚さ計を当てると、底面（内側の面）のこだまの往復の時間は\qty{3.1}{\micro s}でした。厚さはいくらですか。新しいときの厚さが\qty{10}{mm}なら、何 mm 薄くなっていますか。
\par\bkstep{理想化} 配管の壁は平らな板とみる。音の速さは\qty{5900}{m/s}で一定。
\par\kangaekata 往復なので、厚さ＝$ct/2$（④）。
\par\shiki $5900\times3.1\times10^{-6}\div2\approx\qty{9.1e-3}{m}$。$10-9.1=0.9$
\par\bkstep{確かめ} 新しい\qty{10}{mm}なら往復\qty{3.4}{\micro s}。時間が短いので、薄くなっている。
\par\kotae 約\qty{9.1}{mm}、約\qty{0.9}{mm}薄くなっている
\end{reidai}
\begin{jibun}
厚さ\qty{6.0}{mm}の鋼の管の往復の時間は何 μs ですか。はかった時間が\qty{1.8}{\micro s}なら、厚さは何 mm ですか。
\bkjans{$2\times0.0060\div5900\approx\qty{2.0}{\micro s}$。$5900\times1.8\times10^{-6}\div2\approx\qty{5.3e-3}{m}$（約\qty{5.3}{mm}）}
\end{jibun}

\begin{reidai}[ソナーで海の深さをはかる]
\label{ph:波-10:ex7}船のソナーが、真下に出した超音波のこだまを0.20秒後に聞きました。海の深さを求めましょう。
\par\bkstep{理想化} 音はまっすぐ下へ進み、海の底で真上へ返る。海水の音の速さは深さによらず\qty{1500}{m/s}。
\par\kangaekata 往復なので $L=ct/2$（④）。
\par\shiki $L=1500\times0.20\div2$
\par\bkstep{確かめ} 空気（\qty{343}{m/s}）なら同じ時間で\qty{34}{m}。水の中は4倍以上速いので、深くまで測れる。
\par\kotae \qty{150}{m}
\end{reidai}
\begin{jibun}
魚群探知機のこだまが0.060秒後に返ってきました。魚の群れは何 m の深さにいますか。
\bkjans{$1500\times0.060\div2=\qty{45}{m}$}
\end{jibun}

\begin{reidai}[魚の群れの速さ（ドップラー）]
\label{ph:波-10:ex8}魚群探知機が\qty{200}{kHz}の超音波を出したところ、魚の群れからのこだまの振動数が\qty{267}{Hz}高くなっていました。群れは、どちら向きに、何 m/s で動いていますか。
\par\bkstep{理想化} 船は止まっている。群れは、超音波の進む向きにそって動く。動く速さは音の速さよりずっと小さい。
\par\kangaekata 動く物ではね返ると、振動数のずれは $\Delta f\approx2vf/c$（波-8 ⑦）。高くなったのは近づいているから。
\par\shiki $v\approx c\,\Delta f\div(2f)=1500\times267\div(2\times200\times10^3)$
\par\bkstep{確かめ} 約\qty{1.0}{m/s}は、魚が泳ぐ速さとして自然。ずれは\qty{200}{kHz}の約1000分の1.3。
\par\kotae 船に近づく向きに、約\qty{1.0}{m/s}
\end{reidai}
\begin{jibun}
\qty{50}{kHz}の超音波のこだまが、\qty{33}{Hz}低くなりました。群れはどちら向きに、何 m/s で動いていますか。
\bkjans{遠ざかる向き。$1500\times33\div(2\times50000)\approx\qty{0.50}{m/s}$}
\end{jibun}

\begin{reidai}[圧電素子にたまる電気]
\label{ph:波-10:ex9}圧電定数 $d_{33}=\qty{400e-12}{C/N}$ の圧電素子を、\qty{10}{N}の力で押しました。たまる電荷を求めましょう。
\par\bkstep{理想化} 電荷は力に比例する（⑤）。電気はもれない。
\par\kangaekata $Q=d_{33}F$。
\par\shiki $Q=400\times10^{-12}\times10=4.0\times10^{-9}$
\par\bkstep{確かめ} 単位：C/N × N ＝ C。とても小さいので、電子回路で大きくしてはかる。
\par\kotae \qty{4.0e-9}{C}（4.0 nC）
\end{reidai}
\begin{jibun}
水晶（圧電定数 $d_{11}\approx\qty{2.3e-12}{C/N}$）を\qty{10}{N}で押すと、たまる電荷はいくらですか。
\bkjans{$2.3\times10^{-12}\times10=\qty{2.3e-11}{C}$（23 pC。PZT の約170分の1）}
\end{jibun}

\begin{reidai}[ゆれの速度と加速度]
\label{ph:波-10:ex10}機械の床が、\qty{10}{Hz}・振幅\qty{1.0}{mm}の正弦波でゆれています。速度と加速度の振幅を求めましょう。
\par\bkstep{理想化} ゆれは1つの振動数の正弦波とする。
\par\kangaekata $v_{\max}=2\pi fA$、$a_{\max}=(2\pi f)^2A$（⑥）。
\par\shiki $2\pi f=62.8$\,rad/s。$v_{\max}=62.8\times0.0010=0.0628$。$a_{\max}=62.8^2\times0.0010=3.95$
\par\bkstep{確かめ} 加速度は重力の加速度の約0.4倍。手をのせると、はっきりゆれを感じる大きさ。
\par\kotae 速度 約\qty{63}{mm/s}、加速度 約\qty{3.9}{m/s^2}
\end{reidai}
\begin{jibun}
\qty{50}{Hz}・振幅\qty{0.10}{mm}のゆれの、速度と加速度の振幅を求めましょう。
\bkjans{$2\pi\times50\approx314$。$v_{\max}=314\times0.00010\approx\qty{31}{mm/s}$。$a_{\max}=314^2\times0.00010\approx\qty{9.9}{m/s^2}$}
\end{jibun}

\begin{phreidai}[加速度計の感度のくるい]
\label{ph:波-10:ex11}減衰比 $\zeta=0.7$ の加速度計を、固有振動数の0.2倍と0.5倍の振動数で使います。ゆれの倍率 $H$ を求め、感度が何\,\% くるうかを答えましょう。
\phbamen
機械のゆれを加速度計ではかるとき、はかりたい振動数がセンサーの固有振動数に近いと、正しくはかれなくなる。どこまで使えるかを確かめる。
\phmushi
センサーのばね・ダンパーの性質が、振動数で変わらないとする（線形）。ケースは物にしっかりとめてある。
\phzu
振動数の比 $\omega/\omega_{\mathrm n}=0.2,\ 0.5$、減衰比 $\zeta=0.7$、ゆれの倍率 $H$。
\phriyu
加速度計は $|z|=(a_0/\omega_{\mathrm n}^{\,2})H$ で、$H=1$ なら正しい（⑦）。$H$ が1からずれた割合が、感度のくるい。
\phtateru
\[ H=\frac{1}{\sqrt{\bigl(1-(\omega/\omega_{\mathrm n})^2\bigr)^2+\bigl(2\zeta\,\omega/\omega_{\mathrm n}\bigr)^2}} \]
\phtoku
$\omega/\omega_{\mathrm n}=0.2$：$(1-0.04)^2+(2\times0.7\times0.2)^2=0.9216+0.0784=1.0000$ なので $H=1.00$。

$\omega/\omega_{\mathrm n}=0.5$：$(1-0.25)^2+(2\times0.7\times0.5)^2=0.5625+0.49=1.0525$ なので
\[ H=\frac{1}{\sqrt{1.0525}}\approx0.975 \]
\phtashikame
単位なし。大きさ：⑦の「$\zeta=0.7$ では分母が $1+(\omega/\omega_{\mathrm n})^4$」で確かめると、$1+0.5^4=1.0625$、$H\approx0.970$ で、ほぼ同じ。極端な場合：$\omega\to0$ なら $H=1$。
\kotae 0.2倍では $H\approx1.00$（くるいなし）、0.5倍では $H\approx0.97$（約3\,\% 小さく出る）
\end{phreidai}
\begin{jibun}
減衰比 $\zeta=0.1$ の加速度計を、固有振動数の0.2倍と0.5倍で使うと、$H$ はいくらですか。
\bkjans{0.2倍：$0.9216+0.0016=0.9232$、$H\approx1.04$。0.5倍：$0.5625+0.01=0.5725$、$H\approx1.32$（3割も大きく出る）}
\end{jibun}

\begin{reidai}[地震計（変位計）の感度]
\label{ph:波-10:ex12}固有振動数\qty{1.0}{Hz}・減衰比0.7 の変位計（地震計）で、\qty{5.0}{Hz}の地面のゆれをはかります。ずれ $|z|$ は、地面の変位 $Y$ の何倍ですか。
\par\bkstep{理想化} ゆれは1つの振動数の正弦波。センサーは線形。
\par\kangaekata $|z|=(\omega/\omega_{\mathrm n})^2HY$（⑦）。$\omega/\omega_{\mathrm n}=f/f_{\mathrm n}=5$。
\par\shiki $(1-25)^2+(2\times0.7\times5)^2=576+49=625$、$H=1/25$。$(\omega/\omega_{\mathrm n})^2H=25\times\tfrac{1}{25}$
\par\bkstep{確かめ} ちょうど1倍。ばねがやわらかく、おもりがその場に止まっているので、ずれ＝地面の変位。
\par\kotae 1.00倍（地面の変位をそのままはかれる）
\end{reidai}
\begin{jibun}
同じ変位計で、\qty{3.0}{Hz}の地面のゆれをはかると、$|z|$ は $Y$ の何倍ですか。
\bkjans{$(1-9)^2+(1.4\times3)^2=64+17.64=81.64$。$9\div\sqrt{81.64}\approx1.00$（0.996倍）}
\end{jibun}

\begin{reidai}[ポンプのスペクトルを読む]
\label{ph:波-10:ex13}1分に\num{1800}回転するポンプの振動のスペクトルで、\qty{30}{Hz}の線が大きく、\qty{60}{Hz}の線は小さいことがわかりました。何をうたがえばよいですか。
\par\bkstep{理想化} 1倍・2倍の目安（⑧）が当てはまる、よくある原因だけを考える。
\par\kangaekata 回転の周波数 $f_{\mathrm r}=1800\div60$ と、線の振動数を比べる。
\par\shiki $f_{\mathrm r}=\qty{30}{Hz}$。\qty{30}{Hz}は1倍、\qty{60}{Hz}は2倍
\par\bkstep{確かめ} 9-12④の「1倍は重さのかたより、2倍は軸のずれ」と同じ見方。目安なので、ほかの点検と合わせて決める。
\par\kotae 回る部分の重さのかたより（不つり合い）をうたがう
\end{reidai}
\begin{jibun}
1分に\num{1500}回転するモーターのスペクトルで、\qty{50}{Hz}の線が大きくなりました。何倍の成分で、何をうたがいますか。
\bkjans{$f_{\mathrm r}=1500\div60=\qty{25}{Hz}$。\qty{50}{Hz}は2倍なので、軸のつなぎ目のずれ（芯ずれ）}
\end{jibun}

\begin{phreidai}[軸受けの外輪の傷の周波数]
\label{ph:波-10:ex14}玉が9個、玉の直径と玉の中心を結ぶ円の直径の比 $D_{\mathrm b}/D_{\mathrm p}=0.20$ の玉軸受けが、1分に\num{1800}回転の軸を支えています。外輪に傷があると、何 Hz の線が出ますか。内輪の傷なら何 Hz ですか。
\phbamen
ポンプの軸受けがこわれかけていないかを、スペクトルの線の位置で確かめる。
\phmushi
玉のすべり。玉が斜めに押される分（接触角0とする）。
\phzu
玉の数 $n=9$、回転の周波数 $f_{\mathrm r}$、$D_{\mathrm b}/D_{\mathrm p}=0.20$。
\phriyu
外輪の傷は、玉の列（保持器）が回る速さ×玉の数で、たたかれる（⑧）。内輪は、玉の列を追いこしながら回る。
\phtateru
\[ f_{\text{外輪}}=\frac{n}{2}f_{\mathrm r}\Bigl(1-\frac{D_{\mathrm b}}{D_{\mathrm p}}\Bigr),\qquad f_{\text{内輪}}=\frac{n}{2}f_{\mathrm r}\Bigl(1+\frac{D_{\mathrm b}}{D_{\mathrm p}}\Bigr) \]
\phtoku
$f_{\mathrm r}=1800\div60=\qty{30}{Hz}$。
\[ f_{\text{外輪}}=4.5\times30\times0.80=\qty{108}{Hz},\qquad f_{\text{内輪}}=4.5\times30\times1.20=\qty{162}{Hz} \]
\phtashikame
単位：Hz。大きさ：どちらも30の整数倍でないので、不つり合い・芯ずれの線と区別できる。2つをたすと $nf_{\mathrm r}=270$ になる（玉の列と内輪の、たたく回数の合計）。
\kotae 外輪の傷は\qty{108}{Hz}、内輪の傷は\qty{162}{Hz}
\end{phreidai}
\begin{jibun}
玉が8個、$D_{\mathrm b}/D_{\mathrm p}=0.25$ の玉軸受けで、1分に\num{1500}回転の軸を支えています。外輪の傷・内輪の傷の周波数を求めましょう。
\bkjans{$f_{\mathrm r}=\qty{25}{Hz}$。外輪 $4\times25\times0.75=\qty{75}{Hz}$、内輪 $4\times25\times1.25=\qty{125}{Hz}$}
\end{jibun}

\begin{reidai}[ひずみゲージの抵抗の変化]
\label{ph:波-10:ex15}片持ちのはりの根もとに、\qty{120}{\ohm}・ゲージ率2 のひずみゲージをはりました。はりがゆれて、ひずみが $\pm200\times10^{-6}$ の間で変わります。抵抗の変化の振幅を求めましょう。はりが鋼（$E=\qty{205}{GPa}$）なら、応力の振幅はいくらですか。
\par\bkstep{理想化} ゲージは、はりの表面と同じだけ伸び縮みする。ひずみと応力は比例（力-8）。
\par\kangaekata $\Delta R=2\varepsilon R$（⑩）、応力 $=E\varepsilon$。
\par\shiki $\Delta R=2\times200\times10^{-6}\times120=0.048$。応力 $=205\times10^9\times200\times10^{-6}=4.1\times10^7$
\par\bkstep{確かめ} 抵抗の変化は1万分の4。とても小さいので、ブリッジ回路（電-5）で電圧にする。
\par\kotae 抵抗は約\qty{0.048}{\ohm}、応力は約\qty{41}{MPa}
\end{reidai}
\begin{jibun}
\qty{350}{\ohm}・ゲージ率2 のゲージで、ひずみの振幅が $50\times10^{-6}$ のとき、抵抗の変化の振幅は何 Ω ですか。
\bkjans{$2\times50\times10^{-6}\times350=\qty{0.035}{\ohm}$}
\end{jibun}

%% ---------------- まちがえやすいところ ----------------
\Needspace{12\baselineskip}
\bksection{tsumazuki}{まちがえやすいところをチェック}
\begin{tsumazuki}{こだまの距離で、往復を忘れる}
\bkng{鋼で\qty{20}{\micro s}後のこだまなら、傷の深さは $5900\times20\times10^{-6}\approx\qty{118}{mm}$。}
\bkok{行って帰るので、$5900\times20\times10^{-6}\div2\approx\qty{59}{mm}$。}
波-9 と同じまちがいです。ソナー・厚さ計・コウモリでも、いつも2でわります。
\end{tsumazuki}

\begin{tsumazuki}{鋼の中でも、空気の音の速さを使う}
\bkng{鋼の板の厚さ計で、$343\times3.4\times10^{-6}\div2$ を計算する。}
\bkok{音が通っているのは鋼の中なので、鋼の音の速さ\qty{5900}{m/s}を使う。}
距離＝速さ×時間の「速さ」は、音が通っている物の中の速さです。
\end{tsumazuki}

\begin{tsumazuki}{加速度計は、ばねがやわらかいほど感度がよい？}
\bkng{ばねをやわらかくするほど、おもりがよくずれるので、よい加速度計になる。}
\bkok{やわらかいと $\omega_{\mathrm n}$ が下がり、はかりたい振動数が $\omega_{\mathrm n}$ に近づいて、共振でくるう。加速度計は、はかる振動数より $\omega_{\mathrm n}$ をずっと高くする。}
ずれの大きさ（感度）と、使える振動数の広さは、取り引きの関係です。
\end{tsumazuki}

\begin{tsumazuki}{Hz と rad/s をまちがえる}
\bkng{\qty{10}{Hz}・\qty{1}{mm}のゆれの加速度は $10^2\times0.001=\qty{0.1}{m/s^2}$。}
\bkok{$\omega=2\pi f\approx62.8$\,rad/s を使い、$62.8^2\times0.001\approx\qty{3.9}{m/s^2}$。}
$v=A\omega$、$a=A\omega^2$ の $\omega$ は角振動数（rad/s）です。$2\pi$ を忘れると、加速度は約40分の1になってしまいます。
\end{tsumazuki}

%% ---------------- 工学ではこう使う ----------------
\Needspace{8\baselineskip}
\bksection{phgear}{工学ではこう使う}
\begin{phkougaku}[見えない所を見る・ゆれで見守る]
\phboxhead{① 新幹線の車軸の探傷}
新幹線の車軸は、何十万 km も走るうちに、繰り返しの力で小さな割れ（疲労の割れ。力-8）ができることがあります。決まった期間・走行距離ごとの検査で、超音波の探触子を車軸に当てたり、中空の車軸の穴の中に入れたりして、割れがないかを調べています。割れが小さいうちに見つけて、取りかえるのです。

\phboxhead{② 設備をゆれで見守る}
工場や発電所では、大事なポンプや送風機に加速度計をつけっぱなしにして、ゆれのスペクトルをいつも記録しています。いつもと違う線が出てきたり、大きくなってきたりしたら、こわれる前に部品を取りかえます（予知保全）。止まると困る機械ほど、ゆれで「健康診断」をするのです。

\phboxhead{③ 橋とビルのゆれをはかる}
長い橋や高いビルには、加速度計や変位計がとりつけられ、風や地震のときのゆれを記録しています。記録から固有振動数と減衰比を読みとり（⑨）、設計のときの計算と合っているかを確かめます。固有振動数が年とともに下がってきたら、どこかが傷んで、やわらかくなっているかもしれません。
\begin{phketa}[超音波と振動の数の目安]
\phketaline{医用の超音波・鋼の溶接部の探傷の振動数}{約\qtyrange{2}{10}{MHz}・約\qtyrange{2}{5}{MHz}}
\phketaline{鋼の中の\qty{5}{MHz}の波長}{約\qty{1.2}{mm}}
\phketaline{鋼\qty{10}{mm}の往復の時間}{約\qty{3.4}{\micro s}}
\phketaline{圧電式の加速度計で、はかれる上限}{約\qtyrange{5}{50}{kHz}}
\phketaline{毎分1800回転の回転の周波数}{\qty{30}{Hz}}
\phboxhead{答えが桁で合っているか}
鋼の探傷のこだまは、数 μs〜数百 μs。ms（1000分の1秒）の答えになったら、空気の速さを使っていないか、単位の μ を忘れていないかを見直す。機械の振動の加速度が、重力の加速度の何百倍にもなったら、Hz と rad/s を取り違えていないかを見直す。
\end{phketa}
\end{phkougaku}
\phstation{非破壊検査,医用超音波,振動計測,設備診断}{エコーロケーション、パルス反射法、探触子、音響インピーダンス、圧電効果、圧電定数、振動ピックアップ、加速度計、地震計、FFT、不つり合い、芯ずれ、軸受けの欠陥周波数、インパルス応答、打音検査、ひずみゲージ}

\begin{phsaki}{超音波の画像と振動の解析}
この章では、探触子を1つだけ使って、まっすぐ下にある物を見ました。おなかの中の赤ちゃんを動画で見られる超音波の検査では、何十個もの小さな圧電素子を並べ、電気をかけるタイミングを少しずつずらして、超音波の向きを電子的に変えます（フェーズドアレイ）。こだまを向きごとに並べて、体の断面の画像を作るのです。この画像の作り方は、MRI（9-12⑥）と同じく、たくさんの信号から形を組み立てる計算の上に立っています。機械のゆれの解析では、たくさんの点のゆれをいっしょにはかって、ゆれの形（振動モード。波-3）を求めます。こうした方法は、非破壊検査・医用超音波・振動計測の専門の本で学びます。信号をスペクトルに分ける方法は、次の波-11 で学びます。
\end{phsaki}

%% ---------------- 章末問題 ----------------
\bkendproblems
\bkproblemgroup{基本}{3}
\bkproblem コウモリのこだまが4.0 ms 後に返ってきました。虫までの距離を求めましょう（\nobreak\bkreidairef{1}\nobreak）。
\bkpans{$343\times0.0040\div2\approx\qty{0.69}{m}$}
\bkproblem 車の障害物センサーのこだまが1.75 ms 後に返ってきました。壁までの距離を求めましょう（\nobreak\bkreidairef{2}\nobreak）。
\bkpans{$343\times0.00175\div2\approx\qty{0.30}{m}$}
\bkproblem 鋼の板の探傷で、傷のこだまが\qty{12}{\micro s}後に出ました。傷の深さは何 mm ですか（\nobreak\bkreidairef{4}\nobreak）。
\bkpans{$5900\times12\times10^{-6}\div2\approx\qty{0.035}{m}$（約\qty{35}{mm}）}
\bkproblem 魚群探知機のこだまが0.080秒後に返ってきました。魚の群れの深さを求めましょう（\nobreak\bkreidairef{7}\nobreak）。
\bkpans{$1500\times0.080\div2=\qty{60}{m}$}
\bkproblem \qty{25}{Hz}・振幅\qty{0.20}{mm}のゆれの、速度と加速度の振幅を求めましょう（\nobreak\bkreidairef{10}\nobreak）。
\bkpans{$2\pi\times25\approx157$。$157\times0.00020\approx\qty{31}{mm/s}$。$157^2\times0.00020\approx\qty{4.9}{m/s^2}$}
\bkproblem 1分に\num{3000}回転する送風機のスペクトルで、\qty{50}{Hz}の線が大きくなりました。何倍の成分で、何をうたがいますか（\nobreak\bkreidairef{13}\nobreak）。
\bkpans{$f_{\mathrm r}=3000\div60=\qty{50}{Hz}$。1倍なので、回る部分の重さのかたより（不つり合い）}
\bkproblem 圧電定数 $d_{33}=\qty{250e-12}{C/N}$ の圧電素子を\qty{20}{N}で押すと、たまる電荷はいくらですか（\nobreak\bkreidairef{9}\nobreak）。
\bkpans{$250\times10^{-12}\times20=\qty{5.0e-9}{C}$（5.0 nC）}
\bkproblemgroup{チャレンジ}{4}
\bkproblem 水の中に置いた厚さ\qty{20}{mm}の鋼の板に、\qty{30}{mm}上の探触子から超音波を垂直に当てます（水 \qty{1480}{m/s}、鋼 \qty{5900}{m/s}）。(1) 板の表面のこだまと、底面のこだまは、それぞれ何 μs 後に返りますか。(2) 底面のこだまの強さは、出した強さの何\,\% ですか。板の表面を2回通る（入るときと出るとき）とし、底面（鋼→空気）では全部はね返るとします。減衰は無視します（\nobreak\bkreidairef{4}・\bkreidairef{5}\nobreak）。
\bkpans{(1) 表面：$2\times0.030\div1480\approx\qty{40.5}{\micro s}$。底面：$40.5+2\times0.020\div5900$ $\approx40.5+6.8\approx\qty{47.3}{\micro s}$。(2) 1回通るごとに $1-0.88=0.12$。$0.12\times0.12\approx0.014$、約1.4\,\%}
\bkproblem 減衰比 $\zeta=0.7$ の加速度計で、感度のくるいを5\,\% 以内（$H\geqq0.95$）にしたい。$\omega/\omega_{\mathrm n}$ をいくつ以下にすればよいですか。$(\omega/\omega_{\mathrm n})^2=u$ とおいて、2次方程式（3-10）で解きましょう。上限なので、切り捨てて小数第2位まで答えましょう（\nobreak\bkreidairef{11}\nobreak）。
\bkpans{$(1-u)^2+1.96u=1/0.95^2\approx1.108$。$u^2-0.04u-0.108=0$、$u=(0.04+\sqrt{0.0016+0.432})\div2$ $\approx0.349$。$\omega/\omega_{\mathrm n}=\sqrt{0.349}\approx0.591$。0.59 以下}
\bkproblem 玉が7個、$D_{\mathrm b}/D_{\mathrm p}=0.30$ の玉軸受けが、1分に\num{3000}回転の軸を支えています。スペクトルに\qty{227.5}{Hz}の線が出てきました。外輪と内輪のどちらの傷をうたがいますか（\nobreak\bkreidairef{14}\nobreak）。
\bkpans{$f_{\mathrm r}=\qty{50}{Hz}$。外輪 $3.5\times50\times0.70=\qty{122.5}{Hz}$、内輪 $3.5\times50\times1.30=\qty{227.5}{Hz}$。内輪の傷}

\begin{furikaeri}
\begin{checklist}
\item こだまの時間から、コウモリ・ソナー・超音波探傷の距離を求められる。\bkfurref{ph:波-10:sec1}
\item 振動の変位・速度・加速度を換算し、回転機械のスペクトルから故障の見当をつけられる。\bkfurref{ph:波-10:sec6}
\item 1自由度の模型で、加速度計と変位計のしくみを説明し、感度のずれを計算できる。\bkfurref{ph:波-10:sec7}
\end{checklist}
\end{furikaeri}
\bkfukidashi{mina}{見えない所の傷も、機械のゆれも、こだまとセンサーではかれるね。次の波-11 では、はかった音やゆれの信号を、数字にしてスペクトルに分ける方法を学ぶよ。}
